% =====================================================================
% Caffeine Citation: Knowledge Integration in a Cross-Disciplinary Literature
% Specification Outline -- revision 4
% Rewritten for a reader who knows the project and has a seminar's worth of
% network science, and who should not be expected to reconstruct a Bayesian
% multilevel model from notation. Every symbol is defined in prose before it
% appears in an equation. Decisions still open are boxed at the end of the
% section they belong to.
% Causal diagram: caffeine_dag.gv (structure) -> dot2tex -> caffeine_dag.tex.
% =====================================================================
\documentclass[11pt]{article}

\usepackage{graphicx}
\usepackage{float}
\usepackage{pdflscape}
\usepackage{changepage}
\usepackage{ragged2e}
\usepackage{natbib}
\usepackage{titling}
\usepackage{amssymb}
\usepackage{amsmath}
\usepackage{booktabs}
\usepackage{longtable}
\usepackage{underscore}
\usepackage{array}
\usepackage{enumitem}
\usepackage{caption}
\captionsetup{labelfont=bf}
\usepackage{tikz}
\usetikzlibrary{arrows,arrows.meta,shapes,positioning}
\usepackage{bm}
\usepackage[margin=1in]{geometry}

% -- Open items, collected by \listoftodos after the contents --
%    \blocker  : a decision that stops downstream work
%    \checknote: a diagnostic measurement still to be run (C1-C7)
%    \question : an open question or a response not yet incorporated
\usepackage[colorinlistoftodos]{todonotes}
\setlength{\marginparwidth}{0.95in}
\newcommand{\blocker}[1]{\todo[inline, color=red!22, bordercolor=zred]{\textbf{BLOCKING.} #1}}
\newcommand{\checknote}[1]{\todo[color=orange!35, size=\scriptsize]{#1}}
\newcommand{\question}[1]{\todo[color=SkyBlue2!45, size=\scriptsize]{#1}}

\usepackage{fontspec}
\IfFontExistsTF{Latin Modern Roman}{\setmainfont{Latin Modern Roman}}{}

\usepackage[x11names, svgnames, rgb]{xcolor}
\definecolor{gray25}{gray}{0.25}\definecolor{gray45}{gray}{0.45}
\definecolor{gray55}{gray}{0.55}\definecolor{gray60}{gray}{0.60}
\definecolor{gray93}{gray}{0.93}\definecolor{gray98}{gray}{0.98}
\definecolor{zred}{rgb}{0.69,0.19,0.19}
\usepackage{hyperref}
\definecolor{navlink}{rgb}{0.10,0.22,0.55}
\hypersetup{
  colorlinks=true, linkcolor=navlink, citecolor=navlink, urlcolor=navlink,
  linktoc=all, bookmarksnumbered=true,
  pdftitle={Knowledge Integration in the Caffeine Literature: Specification Outline},
  pdfauthor={Jonathan H. Morgan}
}

\newcommand{\customtitle}[1]{%
  \title{\vspace{-1.5em}#1}
  \pretitle{\begin{center}\LARGE\bfseries}
  \posttitle{\par\end{center}\vskip 12pt}
  \preauthor{\begin{center}\normalsize}
  \postauthor{\par\end{center}\vskip 4pt}
  \predate{\begin{center}\small}
  \postdate{\par\end{center}}
}

\newenvironment{springerabstract}{
  \begin{adjustwidth}{2em}{2em}
  \noindent\textbf{Abstract.}\hspace{0.5em}\justifying
}{
  \end{adjustwidth}
}

\newcommand{\springerkeywords}[1]{%
  \begin{adjustwidth}{2em}{2em}
  \vspace{1ex}
  \noindent\textbf{Keywords:} #1
  \end{adjustwidth}
}

\newcommand{\seccap}[1]{{\slshape\small #1\par}\vspace{0.8em}}

% Create Double Rule Column Headers
\newcommand{\doublerule}{%
  \specialrule{\lightrulewidth}{0pt}{\doublerulesep}%
  \specialrule{\lightrulewidth}{0pt}{0pt}}

\newcommand{\meas}{{\footnotesize\textsf{[measured]}}}

% -- Two kinds of open decision, boxed at the end of the section they belong to --
% \begin{decision}  : a substantive judgment the coauthors must make together.
% \begin{measure}   : a measurement-design choice to be settled empirically.
\newcounter{decisioncnt}
\newenvironment{decision}[1]{%
  \refstepcounter{decisioncnt}%
  \par\vspace{0.8em}\noindent
  \begin{adjustwidth}{1em}{1em}%
  \rule{\linewidth}{1.1pt}\par\vspace{0.3em}
  \noindent\textbf{Substantive decision \thedecisioncnt: #1}\par\vspace{0.2em}\small
}{%
  \par\vspace{0.3em}\noindent\rule{\linewidth}{1.1pt}%
  \end{adjustwidth}\vspace{0.8em}
}
\newcounter{measurecnt}
\newenvironment{measure}[1]{%
  \refstepcounter{measurecnt}%
  \par\vspace{0.8em}\noindent
  \begin{adjustwidth}{1em}{1em}%
  \rule{\linewidth}{0.4pt}\par\vspace{0.3em}
  \noindent\textbf{Measurement decision \themeasurecnt: #1}\par\vspace{0.2em}\small
}{%
  \par\vspace{0.3em}\noindent\rule{\linewidth}{0.4pt}%
  \end{adjustwidth}\vspace{0.8em}
}

\newcommand{\code}[1]{\texttt{\small #1}}

\setlength{\droptitle}{0pt}
\customtitle{Knowledge Integration in the Caffeine Literature\\
\large Where Useful Connections Between Research Communities Fail to Form:\\
A Specification Outline}
\author{
  \shortstack{
    Jonathan H. Morgan, Ph.D.$^{1}$ \quad Sarah Delmar, Ph.D.$^{2}$
  }
}
\date{
  $^1$\,Duke Network Analysis Center\\
  $^2$\,Department of Psychology, University of Chicago\\
  \today
}

\begin{document}
\maketitle

\begin{springerabstract}
Caffeine is studied by pharmacologists, food chemists, sleep researchers,
environmental scientists, analytical chemists, and cognitive psychologists. They share
a molecule and very little else. That makes caffeine research a domain without being a
discipline, and it raises a question that a more institutionally coherent field would
not pose so sharply: when one of these communities produces knowledge another could
use, does the knowledge actually travel? Researchers have named specific places where
it should. Ricupero and Ritter, for instance, argue that caffeine's documented effects
on vigilance and appraisal could be implemented in cognitive architectures, and that
the caffeine literature would have to report its results differently for that to
happen. This specification sets out how to test, against the citation record,
whether such connections form. It distinguishes three things that are usually run
together: whether two research communities address related problems, whether they draw
on a common body of prior work, and whether they actually cite one another. Integration can
fail at either transition, and the two failures mean different things.
\end{springerabstract}

\springerkeywords{knowledge integration; citation networks; bibliographic coupling;
community detection; Bayesian multilevel models; caffeine; cognitive science}

\vspace{1.5em}
\begin{adjustwidth}{1em}{1em}
\rule{\linewidth}{1.1pt}\par\vspace{0.4em}
\noindent\textbf{Status note, 22 September 2026 --- start here when picking this work back up.}
\par\vspace{0.3em}\small
The community-detection stage is blocked, and the block is a construction question rather
than a coding one. Two complete CHAMP sweeps over the integration window show that the
Path~A era networks carry no detectable community scale: they are mostly disconnected stars
of one article and the works it cites, the within-era rule discards roughly 70\% of the
citation record, and at the selected resolutions the partition CHAMP returns is the
network's own component structure. The evidence, the three options for rebuilding the
construction, and the seven diagnostic checks (C1--C7) that would decide between them are
all in \S\ref{sec:scalefinding}, at the end of \S\ref{sec:pipeline}. \emph{Nothing
downstream of community detection should be built until that decision is made}, because the
community is the unit of every model here. The tie layer is unaffected: it reads the edge
list and the era maps, not the era networks.
\par\vspace{0.4em}\noindent\rule{\linewidth}{1.1pt}
\end{adjustwidth}

\vspace{1em}
\tableofcontents

\listoftodos[Open Questions and Diagnostic Checks]
\newpage

% =====================================================================
\section{Purpose: Caffeine as a Test Case for Knowledge Integration}
\label{sec:purpose}
\seccap{Why caffeine constitutes a research domain without being a field, and why that
makes integration rather than fragmentation the interesting question.}

\subsection{Caffeine creates a domain, not a discipline}

Caffeine supports many research agendas at once. Its pharmacological profile is unusually
well characterized which makes it useful as a precise experimental instrument. Its
consumption is nearly universal which gives findings about it relevance beyond research
settings. The molecule also appears in enough contexts that very different research
programs encounter it while pursuing very different questions. In the literature,
caffeine is simultaneously a drug, a cognitive enhancer, a food constituent, an
environmental contaminant, and an analytical target.

Caffeine research therefore constitutes a substantive domain without constituting a
single discipline. The distinction is important. Researchers do specialize in caffeine,
journals publish caffeine research repeatedly, and recognizable clusters of investigators
know one another. What the domain lacks is a single institutional structure that
encompasses it as a whole. The questions researchers ask about caffeine, the methods they
use, and the literature they treat as foundational remain strongly shaped by their
parent disciplines.

The result is a domain organized more like a watercooler than a library. Researchers
arrive from different intellectual traditions, exchange knowledge, and move on, without
any structure ensuring that what one group knows reaches another group that could use it.

The venue distribution reflects this structure. Of the 639{,}663 cited works whose
journals can be identified from the Web of Science records, the ten most common journals
account for 7.1\% of the total, and the top fifty account for 19.9\%. Reaching half
requires 374 journals, drawn from 49{,}174 distinct venues \meas{}. Across the full
history of the corpus, the leading venues are broad biomedical and physiological
journals: the \emph{Journal of Biological Chemistry}, \emph{PNAS}, the
\emph{Journal of Physiology}, \emph{Nature}, and \emph{Science}. Within the
2009--2020 window, however, the leading venues shift markedly toward food chemistry,
environmental science, analytical chemistry, and applied biomedical research:
\emph{PLoS ONE}, \emph{Food Chemistry}, \emph{Science of the Total Environment},
the \emph{Journal of Agricultural and Food Chemistry}, the
\emph{Journal of Chromatography~A}, \emph{Water Research}, and
\emph{Chemosphere} \meas{}. A literature organized primarily around one or two research
traditions would not be expected to show either this degree of venue dispersion or this
much change in its leading outlets over time. The evidence does not show that caffeine
research lacks institutions. It shows that no single disciplinary infrastructure
organizes the domain as a whole.

\begin{quote}
\small
A note on what this establishes. The evidence shows that caffeine research is spread
across venues belonging to substantively different domains. It does \emph{not} establish
that caffeine is unusually interdisciplinary, because that is a comparative claim and we
have no matched benchmark. The argument only needs the weaker fact, so we assert only
the weaker fact.
\end{quote}

\subsection{Which makes integration, not fragmentation, the question}

Given all that, finding that the literature is fragmented would be unremarkable. We
should expect distinct communities; the substance does not oblige a pharmacokineticist
and an environmental engineer to read one another.

The interesting question is what happens when they \emph{should}. Sometimes a finding
produced in one corner of this domain is exactly what another corner needs, and the two
corners have every substantive reason to connect. Does the knowledge move?

This is not a bibliometric abstraction we are imposing on the literature. Researchers
have identified the connection points themselves. Ricupero and Ritter's (2024) review argues
that caffeine's measured effects on vigilance, attention, and appraisal could be
implemented directly in cognitive architectures such as ACT-R, which would let system
designers predict how an operator's performance changes over a shift. They also identify
why it has not happened: caffeine studies typically report statistical significance
rather than effect sizes, doses rather than dosages, and no distributional information
about half-life. The knowledge exists; it is not in a form the modeling community can
use. Other reviews in this literature make structurally similar recommendations.

So the questions are: how often are proposed integration points actually realized, where
has integration happened, and what do the successes and the misses have in common?

\subsection{What we need to distinguish, and how bibliometrics lets us}

Answering that requires separating three things that are easy to conflate:

\begin{enumerate}[leftmargin=2em]
  \item whether two research communities \textbf{address related problems};
  \item whether they \textbf{draw on a common body of prior work}; and
  \item whether they \textbf{actually cite and build on one another}.
\end{enumerate}

Two long-standing bibliometric measures map onto the last two.
\textbf{Bibliographic coupling} asks whether two bodies of work cite the same earlier
sources. Two communities that are coupled draw on overlapping bodies of cited literature,
whether or not either is aware of the other. \textbf{Descent} -- direct citation from a
later community to work belonging to an earlier one -- is the stronger relation: a later
community directly references work associated with the earlier one.

The first item, whether the problems are related, has to be measured \emph{outside} the
citation record. Communities are constructed from citation structure, so measuring their
substantive similarity from that same structure would build the answer partly into the
definition. The measure would not be empty -- citation similarity does vary, and
communities differ in density -- but the inference would be circular in exactly the place
where we need it not to be. We measure topical proximity instead from what the articles
say: their abstracts, keywords, and publication venues.

Two further ideas matter later and are worth defining now. A
\textbf{stochastic blockmodel} is a way of estimating how strongly groups connect to one
another while accounting for the fact that large groups have more connections simply by
being large. And a \textbf{generative model} is one that states how the observed counts
could have arisen, which is what lets us say a particular connection is weaker than it
should have been rather than merely small.

\subsection{Caffeine as a case, and the general problem behind it}

Interdisciplinary research is usually justified by the expectation that knowledge
produced in one tradition will become useful in another. A great deal of institutional
effort rests on that expectation, and the outcomes have often fallen short of it.
Existing bibliometric work can count interdisciplinary publications and cross-disciplinary
citations, which tells us where integration \emph{occurred}. It does not tell us where
substantively plausible integration \emph{could have occurred and did not}.

That is the general problem, and caffeine is a case chosen because it offers unusually
favorable conditions for observing it: one shared object, many parent disciplines,
abundant literature, and -- in Ricupero and Ritter and others -- concrete prior claims
that particular connections ought to exist. The method is designed to transfer to other interdisciplinary domains, with the caffeine
case providing its first demonstration.

Finally, this is a literature that changes. Communities appear, grow, merge, and
disappear. We will need to distinguish a \emph{community} in a single year from a
\emph{literature} that persists across years, and from a \emph{lineage} claim that one
year's community is the same community as the next year's. Section~\ref{sec:objects}
does that, and the distinction turns out to determine which questions we can answer with
the data we have.

% =====================================================================
\section{Two Views of the Literature Through Time}
\label{sec:paths}
\seccap{Two ways of building the networks, answering two different kinds of temporal questions.}

The corpus is divided into eras -- for the analytic window, single publication years from
2009 to 2020. How we assign works to eras determines what kind of temporal object we can
study, and there are two defensible constructions that answer different questions. They
are not competing implementations. They are different measurement instruments.

\subsection{Path A: birth-era networks}

A work belongs to the era in which it was published. Each era's network therefore
contains that year's articles and the references entering the literature with them, and
no work appears in two eras.

\checknote{C1: confirm whether cited works are dated by first citation rather than publication. \S\ref{sec:scalefinding}}\emph{This description may not match what the data hold.} Cited works in the node list
appear to carry the year in which they were first cited rather than the year they were
published, which would make an era the year's articles together with the references none of
them had cited before. The evidence is in \S\ref{sec:scalefinding} and the check that would
settle it is C1 there. Until it is settled, read ``published in'' in this subsection as
provisional.

This construction supports questions about the \emph{relationships between} communities
in successive years:

\begin{itemize}[leftmargin=2em]
  \item What knowledge does a later community inherit from an earlier one?
  \item Which communities draw on the same prior works?
  \item When does a shared body of knowledge translate into direct citation?
  \item Where do plausible connections fail to occur?
\end{itemize}

Descent and coupling live here, and so does the whole question of missed opportunities.
This is the path the present analysis takes because the tools needed to construct the
networks and identify their communities already exist.

\subsection{Path B: active-era networks}

A work belongs to \emph{every} era in which it is active -- published in that year or
cited by an article from that year. Works therefore recur across eras, and because they
recur, a community in one year can share members with a community in the next.

That shared membership is what makes it possible to follow a community through time, and
so this construction supports a different class of question:

\begin{itemize}[leftmargin=2em]
  \item Does a research community persist from one year to the next?
  \item Does it grow or contract?
  \item Does it split into several, or merge with others?
  \item Is it absorbed into a neighboring literature?
  \item What predicts how long it survives?
\end{itemize}

\subsection{How the two relate}

The difference is not implementation history. Path A's networks cannot answer Path B's
questions because their communities share no members and therefore have no continuity to
track; Path B's networks cannot sharpen Path A's descent relation because when works
recur across eras, ``cites something that was active in the previous era'' becomes true
of nearly every citation and stops distinguishing anything.

\subsection{Which community partitions the final analysis uses}

The partitions currently on disk were produced at different times by different
procedures -- eras 1--21 under one naming convention, eras 22--23 under another -- and
their community counts swing by a factor of 2.4 across three adjacent eras \meas{}. That
would be a problem for any temporal comparison if those were the partitions we intended
to analyze.

They are not. The final analysis will generate every era's partition under a common
procedure: a Julia implementation of CHAMP, which selects the community-detection
resolution from the data rather than from a value someone typed, and applies the same
selected resolution independently to each era. Because one procedure produces all of
them, community counts become comparable across eras, and a change in the number of
communities becomes a fact about the literature rather than an artifact of method.

The existing partitions remain useful as \emph{development fixtures}: they let the
orchestrator be built and tested against known outputs before the final partitions
exist. They are not the partitions the inference will run on.

\question{Pooled $\gamma^{*}$ cannot be estimated from the Path~A networks. \S\ref{sec:scalefinding}}\emph{The plan described in the rest of this subsection is on hold.} Two CHAMP sweeps over
the integration window show that a pooled $\gamma^{*}$ cannot be estimated from the Path~A
era networks by either of CHAMP's selection criteria, and that the size-invariance argument
below cannot be tested on them. See \S\ref{sec:scalefinding}.

One implementation detail of that plan has to be settled, and it is not a matter of taste.
Under Path~A a work belongs to exactly one era, so the era networks have disjoint vertex
sets and an edge-pooled corpus network is block diagonal. The CHAMP penalty term is computed
against the whole graph's mass, so era $e$ in a pooled sweep sees an effective resolution of
$\gamma\,m_e/m$ rather than $\gamma$ --- roughly a twelfth of the intended value with twelve
eras of comparable size. A $\gamma^{*}$ estimated on the pooled graph is therefore inflated
by about $m/m_e$ and does not transfer back to a single era. We estimate $\gamma^{*}$ instead
from the per-era sweeps and apply the pooled value to every era, which inverts the procedure
of the source manuscript for a reason specific to this construction: there the periods shared
members, so the pooled graph described a real object.

Path B is more than a refactor of the current pipeline: it requires rebuilding the
networks, re-running community detection, and re-generating the community labels, and
the memory demands are large enough that the construction is better done in Julia.
But it is a second analysis of the same corpus, not a separate project. What Path B
invalidates is the \emph{reuse} of Path A's community partitions and their derived files
for Path B's purposes -- not Path A's results, which stand on their own networks.

\begin{quote}
\textbf{\itshape The logical progression across the document runs from the narrowest
relation outward: descent first, then coupling, then -- in Path B -- persistence.}
\end{quote}

% =====================================================================
\section{Objects We Need to Distinguish}
\label{sec:objects}
\seccap{Four things that all get called ``a research area'' in ordinary speech, and why
keeping them apart lets us answer a hard question with easy data.}

Community detection gives us clusters within a single year. Everything else we might want
to say about how research areas behave over time is a further claim, and the claims
differ in how much they cost.

\begin{center}
\small
\begin{tabular}{@{}lp{0.70\linewidth}@{}}
\toprule
\textbf{Object} & \textbf{What it is, and what claiming it commits us to}\\
\doublerule
Article    & An observable document: a caffeine paper, or a work one cites. No
             inference involved.\\[3pt]
Community  & A set of articles clustered together \emph{within one era}. Under Path A,
             calling something ``Community 17'' in 2016 makes no claim whatever that
             there is a Community 17 in 2017. The label is local to the year.\\[3pt]
Literature & A higher-order grouping of communities that occupy a common position in
             the citation structure. Several communities across different years may all
             belong to a pharmacokinetic literature without any one of them persisting
             intact. Identified by a nested stochastic blockmodel
             (\S\ref{sec:literatures}), \emph{not} by journal membership.\\[3pt]
Lineage    & The stronger claim that community $A$ in 2016 \emph{is} the predecessor of
             community $B$ in 2017. This requires a rule for establishing identity through
             time, which Path A cannot supply.\\
\bottomrule
\end{tabular}
\end{center}

The definition of a literature deserves one further sentence, because it is easy to
assume we are simply grouping journals. We are not. Literatures are found from the
citation structure itself, using a blockmodel that groups communities occupying similar
positions in it. Journal composition and community labels then serve to \emph{interpret}
those groupings and to test whether they correspond to recognizable disciplinary
boundaries. Discovery and validation stay separate, and the validation is meaningful only
because the journal information played no part in the discovery.

The payoff is worth stating plainly, because it determines the shape of the project:

\begin{quote}
We can study the rise and decline of a \emph{literature} without solving the much harder
problem of establishing \emph{community} identity through time.
\end{quote}

A literature's lifecycle is measured by how much of the field it occupies in each year --
how many articles and how many communities belong to it, and what share of citation
traffic passes through it. Those quantities can rise and fall without our ever claiming
that any particular community survived. Community-level survival, which does require the
identity claim, is Path B's business (\S\ref{sec:pathb}).

% =====================================================================
\section{Research Questions}
\label{sec:questions}
\seccap{Why a baseline has to come before any claim that a connection is missing.}

\subsection{Two temporal windows}

The corpus spans 23 eras, from 1907--1944 through 2020. Those eras are not equal in
width: the early ones cover decades, and eras 12--23 are the only unbroken run of single
years. Two different
spans are therefore in play, and confusing them makes several questions ambiguous.

\begin{center}
\small
\begin{tabular}{@{}lp{0.66\linewidth}@{}}
\toprule
\textbf{Window} & \textbf{What it covers and what it is used for}\\
\doublerule
Historical window & All 23 eras. Used to reconstruct the major literatures and describe
                    how they emerge, grow, contract, and decline over the full history of
                    the field.\\[3pt]
Integration window & Eras 12--23, single years from 2009 to 2020, giving eleven adjacent
                    era pairs on a constant one-year base. Used for the models of
                    integration, where comparing across pairs requires that the pairs be
                    commensurable.\\
\bottomrule
\end{tabular}
\end{center}

Era width matters for the integration models because the descent relation looks one era
back: a one-year window and a thirty-eight-year window select very different citations.
It matters much less for describing a literature's trajectory, where the quantity is how
much of the field the literature occupies and the eras are simply time points.

\subsection{Why the questions come in tiers}

We cannot call a connection ``missing'' merely because two communities fail to cite one
another. Most communities probably should not cite one another -- an environmental
chemistry cluster and a sleep-deprivation cluster have little to say to each other, and
their mutual silence is not a failure of anything.

So a claim about a missed opportunity requires three things in order. First, we need to
know how integration \emph{normally} works in this literature: given two communities with
comparable overlap in what they read, how much citation ordinarily follows? Second, we
need to know where those communities sit in the larger intellectual structure, because a
shortfall between two clusters of the same literature means something different from a
shortfall across a disciplinary boundary. Only then can we ask whether cognitive research
departs from the pattern.

That is why the questions come in three tiers, and why the third is uninterpretable
without the first.

\subsection{Tier 1: what normal integration looks like}

\begin{itemize}[leftmargin=2em]
  \item \textbf{RQ1.} Do communities that work on related problems end up drawing on a
        common body of prior work, and how strongly?
  \item \textbf{RQ2.} Given a common body of prior work, how much direct citation
        ordinarily follows? \emph{This is the baseline every later claim is read
        against.}
  \item \textbf{RQ3.} Are communities that draw heavily on the past also the ones the
        future draws on? That is, are being derivative and being generative related
        traits?
  \item \textbf{RQ4.} Does either conversion strengthen or weaken across the integration
        window?
\end{itemize}

\subsection{Tier 2: the intellectual structure the communities sit in}

\begin{itemize}[leftmargin=2em]
  \item \textbf{RQ5.} What are the major literatures in caffeine research, and how do
        they emerge, grow, contract, and decline across the historical window -- the full
        23 eras, not only the recent ones?
  \item \textbf{RQ6.} Is the pharmacokinetic literature as large and persistent as the
        argument assumes? This tests the premise of the project rather than its
        conclusion -- if that literature turns out to be small or transient, the claim
        that it offers underused knowledge weakens before cognitive science is examined.
  \item \textbf{RQ7.} Does the problems-to-shared-reading conversion differ among
        literatures? And do the communities the citation data produce line up with the
        disciplinary boundaries that journals imply, or cut across them?
\end{itemize}

\subsection{Tier 3: whether cognitive research is different}

\begin{itemize}[leftmargin=2em]
  \item \textbf{RQ8.} Where does cognitive research sit -- gathered into its own
        communities, or scattered through communities organized around other concerns?
  \item \textbf{RQ9.} Is the conversion from shared knowledge into citation
        \emph{asymmetric} where cognitive research is involved? \emph{This is the primary
        claim, and unlike the others it comes with a stated expectation.}
  \item \textbf{RQ10.} Which specific pairs of communities show a connection that should
        have formed and did not, and is cognitive research overrepresented among them?
  \item \textbf{RQ11.} Has the shortfall changed across the integration window?
\end{itemize}

\subsection{What we expect to find, and why}

RQ9 is not a neutral question about which side shows a shortfall. We expect a particular
asymmetry, and stating it in advance is what makes the finding informative either way.

\begin{quote}
\textbf{Expectation.} Conditional on shared knowledge and comparable opportunity,
cognitive communities should be more likely to cite relevant neuroscientific and
pharmacological caffeine research than those literatures are to incorporate higher-level
cognitive constructs.
\end{quote}

The reason is a difference in the level of abstraction at which the two literatures
operate. Neuroscience and pharmacology characterize lower-level mechanisms -- receptor
binding, adenosine antagonism, coordination among neural systems. Cognitive science
organizes such mechanisms under higher-order constructs: working memory, sustained
attention, appraisal. A cognitive researcher has a standing reason to reach downward for
mechanistic evidence, because the mechanism is what the construct is supposed to be
implemented in. A neuroscientist studying receptor kinetics does not need to reach upward
to a cognitive architecture in order to answer the question in front of them.

If that is right, the shortfall should appear on one side and not the other, and the
direction is predicted rather than discovered. A symmetric shortfall, or one running the
other way, would be evidence against this account of how the two fields relate.

\subsection{One limit the corpus imposes on RQ9}

The corpus consists of caffeine papers plus everything those papers cite. Work that never
mentions caffeine is not in it at all. That asymmetry constrains what RQ9 can observe, and
the constraint is worth making concrete.

Suppose cognitive caffeine researchers cite pharmacokinetic caffeine research. We can see
that, because both sides are inside a caffeine-centered corpus.

Now suppose the reverse: pharmacokinetic researchers should be incorporating a general
cognitive model -- a theory of vigilance, say -- that never mentions caffeine. That model
is outside the frame entirely. Its absence from the corpus is guaranteed by how the corpus
was built and tells us nothing about whether anyone should have cited it.

So absence cannot be read symmetrically. We can measure whether cognitive communities cite
caffeine work, and whether caffeine communities cite cognitive work \emph{that is
already inside the caffeine literature}. We cannot measure whether they should have reached
outside it. Any claim about direction is reported as a bound, not an estimate.

\question{Sarah responded 22 September 2026; her answer is not yet written into the document.}
\begin{decision}{Whether the predicted asymmetry is theoretically right}
The expectation above rests on a characterization of how cognitive science and
neuroscience relate -- that one organizes the other's mechanisms under higher-order
constructs, and that the reaching-down relation is more natural than the reaching-up one.
That is a claim about the two fields, not about the data, and Sarah is far better placed
than the citation record to judge whether it is accurate. If the characterization is
wrong the model still estimates both directions, but the paper would be reporting a
discovery rather than testing a prediction, which is a weaker result.
\end{decision}

% =====================================================================
\section{A Conceptual Model of Knowledge Integration}
\label{sec:causal}
\seccap{The argument as a diagram: what we think produces a citation between two
communities, and which parts of that we can observe.}

\begin{landscape}
\begin{figure}[p]
\centering
\begin{tikzpicture}[>=latex',line join=bevel,]
%%
\begin{scope}
  \pgfsetstrokecolor{black}
  \definecolor{strokecol}{rgb}{1.0,1.0,1.0}
  \pgfsetstrokecolor{strokecol}
  \definecolor{fillcol}{rgb}{1.0,1.0,1.0}
  \pgfsetfillcolor{fillcol}
  \fill [opacity=1] (0.0bp,0.0bp) -- (0.0bp,344.82bp) -- (644.32bp,344.82bp) -- (644.32bp,0.0bp) -- cycle;
  \draw [opacity=0.0] (0.0bp,0.0bp) -- (0.0bp,344.82bp) -- (644.32bp,344.82bp) -- (644.32bp,0.0bp) -- cycle;
\end{scope}
\begin{scope}
  \pgfsetstrokecolor{black}
  \definecolor{strokecol}{rgb}{1.0,1.0,1.0}
  \pgfsetstrokecolor{strokecol}
  \definecolor{fillcol}{rgb}{1.0,1.0,1.0}
  \pgfsetfillcolor{fillcol}
  \fill [opacity=1] (0.0bp,0.0bp) -- (0.0bp,344.82bp) -- (644.32bp,344.82bp) -- (644.32bp,0.0bp) -- cycle;
  \draw [opacity=0.0] (0.0bp,0.0bp) -- (0.0bp,344.82bp) -- (644.32bp,344.82bp) -- (644.32bp,0.0bp) -- cycle;
\end{scope}
\begin{scope}
  \pgfsetstrokecolor{black}
  \definecolor{strokecol}{rgb}{1.0,1.0,1.0}
  \pgfsetstrokecolor{strokecol}
  \definecolor{fillcol}{rgb}{1.0,1.0,1.0}
  \pgfsetfillcolor{fillcol}
  \fill [opacity=1] (0.0bp,0.0bp) -- (0.0bp,344.82bp) -- (644.32bp,344.82bp) -- (644.32bp,0.0bp) -- cycle;
  \draw [opacity=0.0] (0.0bp,0.0bp) -- (0.0bp,344.82bp) -- (644.32bp,344.82bp) -- (644.32bp,0.0bp) -- cycle;
\end{scope}
\begin{scope}
  \pgfsetstrokecolor{black}
  \definecolor{strokecol}{rgb}{1.0,1.0,1.0}
  \pgfsetstrokecolor{strokecol}
  \definecolor{fillcol}{rgb}{1.0,1.0,1.0}
  \pgfsetfillcolor{fillcol}
  \fill [opacity=1] (0.0bp,0.0bp) -- (0.0bp,344.82bp) -- (644.32bp,344.82bp) -- (644.32bp,0.0bp) -- cycle;
  \draw [opacity=0.0] (0.0bp,0.0bp) -- (0.0bp,344.82bp) -- (644.32bp,344.82bp) -- (644.32bp,0.0bp) -- cycle;
\end{scope}
\begin{scope}
  \pgfsetstrokecolor{black}
  \definecolor{strokecol}{rgb}{1.0,1.0,1.0}
  \pgfsetstrokecolor{strokecol}
  \definecolor{fillcol}{rgb}{1.0,1.0,1.0}
  \pgfsetfillcolor{fillcol}
  \fill [opacity=1] (0.0bp,0.0bp) -- (0.0bp,344.82bp) -- (644.32bp,344.82bp) -- (644.32bp,0.0bp) -- cycle;
  \draw [opacity=0.0] (0.0bp,0.0bp) -- (0.0bp,344.82bp) -- (644.32bp,344.82bp) -- (644.32bp,0.0bp) -- cycle;
\end{scope}
  \definecolor{fillcolor}{rgb}{0.89,0.85,0.94}
  \node (Cog) at (473.0bp,319.37bp) [draw,fill=fillcolor,ellipse] {\shortstack{Cognitive\\[-1pt]orientation}};
  \node (V) at (435.0bp,217.46bp) [draw,fill=gray93,ellipse] {$V_{ij}$ venue};
  \node (U) at (578.0bp,217.46bp) [draw,fill=gray93,ellipse] {$U_{ij}$ authors};
  \definecolor{strokecolor}{rgb}{0.72,0.53,0.04}
  \definecolor{fillcolor}{rgb}{0.97,0.91,0.74}
  \node (S) at (284.0bp,217.46bp) [draw=strokecolor,fill=fillcolor,ellipse] {\shortstack{Topical proximity\\[-1pt]$S_{ij}$}};
  \node (T) at (248.0bp,113.0bp) [draw,fill=gray93,ellipse] {Era $p$};
  \definecolor{strokecolor}{rgb}{0.48,0.12,0.12}
  \node (A) at (123.0bp,113.0bp) [draw=strokecolor,ellipse,dashed,fill=white] {\shortstack{$A_{ij}$\\[-1pt]{\scriptsize\itshape latent}}};
  \definecolor{strokecolor}{rgb}{0.18,0.49,0.2}
  \definecolor{fillcolor}{rgb}{0.85,0.92,0.83}
  \node (EK) at (27.0bp,113.0bp) [draw=strokecolor,fill=fillcolor,ellipse] {$E^{K}_{ij}$};
  \definecolor{strokecolor}{rgb}{0.18,0.49,0.2}
  \definecolor{fillcolor}{rgb}{0.85,0.92,0.83}
  \node (EC) at (332.0bp,113.0bp) [draw=strokecolor,fill=fillcolor,ellipse] {$E^{C}_{ij}$};
  \node (W) at (452.0bp,113.0bp) [draw,fill=gray93,ellipse] {$u_h$};
  \node (K) at (182.0bp,22.0bp) [draw,fill=white,ellipse,double,double distance=1.2pt,line width=0.9pt] {\shortstack{Coupling\\[-1pt]$K_{ij}$}};
  \node (C) at (399.0bp,22.0bp) [draw,fill=gray98,ellipse,double,double distance=1.2pt,line width=0.9pt] {\shortstack{Descent\\[-1pt]$C_{hij}$}};
  \definecolor{strokecolor}{rgb}{0.42,0.25,0.63}
  \draw [strokecolor,->] (Cog) ..controls (404.78bp,282.31bp) and (360.87bp,259.09bp)  .. (S);
  \definecolor{strokecolor}{rgb}{0.42,0.25,0.63}
  \draw [strokecolor,->] (Cog) ..controls (567.5bp,298.55bp) and (613.58bp,279.46bp)  .. (637.0bp,242.91bp) .. controls (649.21bp,223.86bp) and (643.35bp,213.72bp)  .. (637.0bp,192.0bp) .. controls (617.47bp,125.18bp) and (601.84bp,100.97bp)  .. (540.0bp,69.0bp) .. controls (513.73bp,55.422bp) and (482.66bp,44.769bp)  .. (C);
  \draw [gray55,->] (V) ..controls (337.04bp,184.29bp) and (226.56bp,148.01bp)  .. (A);
  % hand-placed: V -> C, direct venue path (see caffeine_dag.tex, revision 4)
  \draw [gray55,->] (V) ..controls (409.0bp,167.0bp) and (392.0bp,88.0bp)  .. (C);
  \draw [gray55,->] (U) ..controls (524.13bp,200.04bp) and (508.44bp,195.61bp)  .. (494.0bp,192.0bp) .. controls (365.8bp,159.95bp) and (332.3bp,158.29bp)  .. (203.0bp,131.0bp) .. controls (192.58bp,128.8bp) and (181.39bp,126.43bp)  .. (A);
  \draw [gray55,->] (U) ..controls (575.38bp,165.54bp) and (568.97bp,123.8bp)  .. (548.0bp,95.0bp) .. controls (526.39bp,65.322bp) and (489.23bp,47.332bp)  .. (C);
  \draw [gray60,->] (T) ..controls (202.97bp,126.48bp) and (185.66bp,127.54bp)  .. (A);
  \draw [gray60,->] (T) ..controls (292.89bp,85.542bp) and (331.36bp,62.865bp)  .. (C);
  \definecolor{strokecolor}{rgb}{0.72,0.53,0.04}
  \draw [strokecolor,->] (S) ..controls (224.33bp,178.49bp) and (183.06bp,152.22bp)  .. (A);
  \definecolor{strokecol}{rgb}{0.0,0.0,0.0}
  \pgfsetstrokecolor{strokecol}
  \draw (222.5bp,161.5bp) node {$\bm{\alpha}_S$};
  \definecolor{strokecolor}{rgb}{0.72,0.53,0.04}
  \draw [strokecolor,->] (S) ..controls (390.95bp,182.41bp) and (497.38bp,146.43bp)  .. (508.0bp,131.0bp) .. controls (531.78bp,96.472bp) and (485.13bp,63.579bp)  .. (C);
  \draw (528.0bp,113.0bp) node {{\scriptsize direct}};
  \definecolor{strokecolor}{rgb}{0.48,0.12,0.12}
  \draw [strokecolor,->] (A) ..controls (142.17bp,83.084bp) and (153.3bp,66.291bp)  .. (K);
  \definecolor{strokecolor}{rgb}{0.48,0.12,0.12}
  \draw [strokecolor,->] (A) ..controls (203.5bp,86.043bp) and (290.35bp,58.036bp)  .. (C);
  \definecolor{strokecolor}{rgb}{0.18,0.49,0.2}
  \draw [strokecolor,->] (EK) ..controls (70.909bp,86.788bp) and (111.69bp,63.372bp)  .. (K);
  \definecolor{strokecolor}{rgb}{0.18,0.49,0.2}
  \draw [strokecolor,->] (EC) ..controls (352.95bp,84.174bp) and (366.18bp,66.596bp)  .. (C);
  \draw [gray60,->] (W) ..controls (372.93bp,85.936bp) and (289.29bp,58.367bp)  .. (K);
  \draw [gray60,->] (W) ..controls (434.9bp,83.281bp) and (425.1bp,66.834bp)  .. (C);
%
%
  % -- hand-drawn bow arc: residual common causes beyond A and U --
  \draw[<->,dashed,zred,line width=1.0pt] (K) to[bend right=28]
       node[pos=0.5,below=1pt,text=black,font=\small]
       {residual $Z$} (C);
\end{tikzpicture}
\caption{The conceptual model. Topical proximity $S_{ij}$ shapes a shared knowledge base
$A_{ij}$, which we cannot observe (dashed). Coupling $K_{ij}$ and descent $C_{hij}$ are
two observed manifestations of it (double outlines). Remaining arrows are the alternative
reasons a connection might form: shared authors $U_{ij}$, shared venues $V_{ij}$,
opportunity $E^{K}$ and $E^{C}$, individual work salience $u_h$, and era $p$. Drawn in
Pearl's semi-Markovian notation; the dashed double-headed arc marks common causes of both
indicators that we cannot measure, such as reciprocal reviewing obligations and shared
institutional context.}
\label{fig:dag}
\end{figure}
\end{landscape}

\subsection{Reading the causal diagram from left to right}

The diagram organizes the argument around a latent \textbf{shared knowledge base}
$A_{ij}$ and its two observable manifestations, \textbf{coupling} $K_{ij}$ and
\textbf{descent} $C_{hij}$. The variables above and around them represent the
conditions that can shape those relationships.

Start with \textbf{topical proximity} $S_{ij}$, shown in the upper-left portion of
the diagram. Communities that study similar problems have more reason to encounter
the same underlying work. We measure that proximity from what their articles say
rather than from whom they cite. Topical proximity can affect integration in two
ways. First, it can shape the shared knowledge base $A_{ij}$. Second, the diagram
allows a direct relationship with descent, recognizing that communities working on
similar problems may cite one another for reasons not fully captured by their shared
references.

The \textbf{shared knowledge base} $A_{ij}$ appears below topical proximity. It
represents the overlap in what the two communities actually know about. We cannot
observe that knowledge directly; nobody records everything researchers have read or
know. The dashed outline therefore marks $A_{ij}$ as latent.

At the bottom of the diagram are the two relationships we can observe.
\textbf{Bibliographic coupling} $K_{ij}$ provides one trace of the shared knowledge
base. When two communities cite the same earlier works, they reveal common
intellectual foundations. \textbf{Descent} $C_{hij}$ provides a second, stronger
form of integration: a later community directly cites work belonging to an earlier
community. The later community is therefore not merely drawing on some of the same
background literature; it is citing the earlier community's output directly.

The remaining variables represent other reasons these observed relationships may
form. \textbf{Cognitive orientation}, shown at the top of the diagram, can shape
both the problems a community studies and its propensity for direct citation.
Shared authors $U_{ij}$ and shared venues $V_{ij}$ provide additional routes to
connection. Opportunity $E^{K}_{ij}$ and $E^{C}_{ij}$ records how many chances the
two observed relationships had to form. Work salience $u_h$ allows some individual
papers to attract citations across communities regardless of their substantive
relationship, while era $p$ accounts for temporal differences in the opportunity
to accumulate citations.

Finally, the dashed double-headed arc between coupling and descent represents
unmeasured common causes of both relationships. These may include shared
institutional settings, reciprocal reviewing obligations, or other forms of
intellectual contact that the citation record does not capture. The arc is important
because it reinforces the central causal distinction developed in the next section:
coupling and descent may covary without coupling itself causing descent.

\subsection{Why coupling is evidence, not a cause}

One feature of the diagram matters enough to state separately, because it changes what we
are allowed to claim.

It is tempting to say that shared reading \emph{causes} citation -- that $K$ leads to $C$.
But bibliographic coupling is not something that happens to two communities before they
cite each other. It is an observed consequence of their citation behavior, computed after
the fact from the same reference lists. In the diagram, $K$ and $C$ are two manifestations
of the same underlying thing, $A$. They covary because they share a cause, not because one
produces the other.

\begin{quote}
The shared knowledge base is a theoretical construct, not an estimated score. Coupling is
\emph{evidence} that it exists. We do not claim that coupling itself causes citation.
\end{quote}

This has a practical consequence: $A$ is not fitted. Two indicators are not enough to
identify a latent variable without additional assumptions we have no grounds for, so we
use $K$ as an observable stand-in and report the relationship between $K$ and $C$ as an
association -- how much citation accompanies a given degree of shared reading -- rather
than as a causal effect. Everything the argument needs survives that weakening.

\subsection{Where the cognitive hypothesis lives}

A causal diagram shows which variables affect which. It is not well suited to showing the
claim this project actually makes, which is about the \emph{strength} of a relationship
rather than its existence. Figure~\ref{fig:moderation} states that separately.

\begin{figure}[H]
\centering
\begin{tikzpicture}[
  var/.style={draw, rounded corners=2pt, inner sep=5pt, font=\small, align=center},
  ind/.style={var, double, double distance=1.2pt},
  mod/.style={var, fill=black!5},
  arr/.style={-{Stealth[length=2.4mm]}, thick}]
  \node[ind] (k) at (0,0) {Coupling $K_{ij}$\\{\scriptsize shared reading}};
  \node[ind] (c) at (7,0) {Descent $C_{hij}$\\{\scriptsize citation}};
  \node[mod] (m) at (3.5,2.0) {Cognitive orientation of the two communities};
  \draw[arr] (k) -- node[below, font=\footnotesize] {how much citation follows} (c);
  \draw[arr, zred] (m) -- (3.5,0.42);
  \node[right=2pt, font=\footnotesize, text=zred] at (3.6,1.1) {weakens this?};
\end{tikzpicture}
\caption{The claim under test is not that cognitively-oriented communities cite less in
general. It is that the \emph{conversion} of shared reading into citation is weaker where
cognitive research is involved -- they read the same works, and citation does not follow.
Statistically this is an interaction, which is why it cannot be drawn as an arrow in
Figure~\ref{fig:dag}.}
\label{fig:moderation}
\end{figure}



\subsection{Two places where integration can fail}
\label{sec:twoplaces}

The chain has two transitions, and a shortfall at each means something different. This
distinction is the analytical core of the project.

\paragraph{Stage 1 -- from related problems to shared knowledge.}
If two communities work on similar problems but read largely different literatures, the
barrier arises \emph{before} any question of direct citation. They are not positioned to
integrate because they do not share foundations. Nothing is being ignored; the two
research programs have grown from different roots.

\paragraph{Stage 2 -- from shared knowledge to direct citation.}
If two communities already draw on many of the same works and nevertheless do not cite one
another, the situation is different. They have common intellectual foundations, and those
foundations are not producing integration. This is the pattern that deserves to be called
a missed opportunity, and it is what Ricupero and Ritter describe: a body of caffeine
research whose findings the modeling community could use, built on overlapping
foundations, yet not directly cited.

Naming these separately is what makes the finding interpretable either way it comes out.

\subsection{Alternative explanations we have to rule out}

Before attributing a shortfall to non-integration, four other mechanisms have to be
considered. Each predicts something we can look for.

\begin{center}
\small
\begin{tabular}{@{}llp{0.50\linewidth}@{}}
\toprule
 & \textbf{Mechanism} & \textbf{What it predicts}\\
\doublerule
M1 & Vintage &
  Cognitive work in this corpus is more recent, so less of it has been available long
  enough to accumulate citations. If so, the shortfall shrinks once we account for how
  old the cited work is.\\[3pt]
M2 & Separation &
  High topical proximity, unexpectedly low coupling. The communities work on related
  problems and read different literatures. This is Stage~1 failing.\\[3pt]
M3 & Non-integration &
  High coupling, unexpectedly low citation. They hold foundations in common and citation
  does not follow. This is Stage~2 failing.\\[3pt]
M4 & Opportunity &
  Coupling and citation are exactly what the communities' sizes and activity levels
  predict. Nothing needs explaining. This is the null.\\[3pt]
M5 & Disciplinary distance &
  The shortfall tracks how far apart two communities are \emph{disciplinarily}, not
  whether either is cognitive. If so, the cognitive result is an instance of a general
  norm rather than a fact about cognitive science.\\
\bottomrule
\end{tabular}
\end{center}

M2 and M3 are the substantively interesting outcomes, and they are separated by
\emph{which} transition falls short -- which is why the two stages are estimated
separately rather than collapsed into one relationship. M4 is why Tier~1 has to come
first. M5 is a genuine rival to the cognitive claim, and testing it requires disciplinary
distance in the model alongside cognitive orientation.

These are mechanisms, not formal hypotheses. The formal questions are the RQs of
\S\ref{sec:questions}; M1--M5 are the competing stories those questions have to
adjudicate between.

% =====================================================================
\section{Constructing the Networks and Measures}
\label{sec:networks}
\seccap{What exactly gets built: the nodes, the ties, and how each measure is computed.}

\subsection{The two network types}

Three things are easy to run together here, so it is worth being explicit about the
counting. There are \textbf{two network types} -- descent and coupling -- distinguished by
what makes a tie. There are \textbf{three tie weightings} -- breadth, fractional, and
intensity -- distinguished by what a tie's value counts. Each weighting applied to each
type produces a different weighted network, so

\[
2 \text{ network types} \times 3 \text{ weightings} = 6 \text{ weighted networks per era
pair.}
\]

All six share the same nodes. Within a type, all three weightings share the same ties as
well, so one layout and one role partition serve all three.

For each adjacent pair of eras $(t{-}1, t)$, Path A produces two community-level networks
over the same conceptual population.

\paragraph{The descent network.}
\begin{itemize}[leftmargin=2em, itemsep=1pt]
  \item \textbf{Nodes.} Communities from the two eras. A community appearing in both roles
        is two nodes, one per era.
  \item \textbf{Tie.} An era-$t$ community cites a work that belonged to an era-$(t{-}1)$
        community and was actively cited during that earlier era.
  \item \textbf{Direction.} Directed, from the citing community to the cited one, following
        the citation.
  \item \textbf{Weight.} Yes -- see \S\ref{sec:weightings}.
  \item \textbf{Meaning.} Direct intellectual uptake. The later community reached into
        material that was part of the earlier community's own structure.
\end{itemize}

\paragraph{The coupling network.}
\begin{itemize}[leftmargin=2em, itemsep=1pt]
  \item \textbf{Nodes.} The same communities.
  \item \textbf{Tie.} The two communities' articles cite at least one work in common,
        wherever that work sits.
  \item \textbf{Direction.} \emph{Coupling is mathematically symmetric.} Sharing a
        reference is a property of the pair, not of one member. The Path A output orients
        it from the later era to the earlier one solely so that it aligns with the descent
        network and the two can be compared tie by tie; that orientation is a
        visualization convention and does not make coupling a directed relation. Because
        it is symmetric, a weight normalized within the later community and one normalized
        within the earlier community are both meaningful, and both are reported.
  \item \textbf{Weight.} Yes.
  \item \textbf{Meaning.} Common intellectual foundations, with no requirement that either
        community cites the other.
\end{itemize}

\paragraph{How restrictive descent is.} Of the citations an era emits, 0.54\% and 1.55\%
survive into the descent network on the two era pairs examined so far \meas{}. That is the
price of the strong definition: the cited work must both belong to the earlier era and have
been in active use there. It is an indication of how narrow a one-year inheritance window
is, and not a clean measurement of citation lag -- it also reflects how the corpus was
built and how the eras were partitioned, and separating those would take a decomposition
we have not done.

\subsection{Three ways to weight each network}
\label{sec:weightings}

Each tie can be weighted three ways, and each answers a different question. All three use
the same nodes and the same ties; only the numbers change.

\begin{center}
\small
\begin{tabular}{@{}lp{0.38\linewidth}p{0.34\linewidth}@{}}
\toprule
\textbf{Weighting} & \textbf{What is counted} & \textbf{Question it answers}\\
\doublerule
Breadth    & Distinct shared works, each counted once for every community that cites it &
             How widely does one community draw on the other?\\[3pt]
Fractional & Each shared work is one unit of evidence, divided equally among the community
             pairs it connects &
             How was one community's contribution distributed among its heirs?\\[3pt]
Intensity  & Citation events rather than distinct works; for coupling, the product of the
             number of citing articles on each side &
             How heavily does one community lean on the other?\\
\bottomrule
\end{tabular}
\end{center}

Fractional weighting has one property the others lack: its totals do not change when the
community partition is finer or coarser, because each shared work contributes exactly one
unit however many pairs it spans. Breadth totals grow with the number of communities, and
community counts vary by a factor of 2.4 across three adjacent eras \meas{}, so breadth
totals are not comparable across pairs. Where a cross-year comparison is wanted, fractional
is the weighting that supports it.

\subsection{Topical proximity, measured outside the citation record}
\label{sec:proximity}

Topical proximity $S_{ij}$ is how similar two communities' subject matter is, computed from
the articles' abstracts, keywords, and journals -- never from their references.

That independence is the whole point. Communities are \emph{detected} from the citation
network, so measuring their similarity from citations too would build the answer partly
into the definition. Measuring similarity from text gives an independent yardstick, and
only then does ``these two communities are substantively close and do not connect'' become
a claim with content.

Naming the data source is not the same as specifying the measure, though, and this measure
carries more weight than any other single quantity in the design: $S_{ij}$ is what
separates mechanism M2 from M3, and it is the covariate that makes a shortfall
interpretable. It deserves to be built deliberately.

\paragraph{The problem a naive measure runs into.}
Every community in this corpus talks about caffeine. A conventional similarity measure over
term-frequency vectors will therefore find every pair somewhat similar, because they share
the domain's ubiquitous vocabulary -- \emph{caffeine}, \emph{coffee}, \emph{consumption},
\emph{dose}. Worse, it rewards globally frequent terms, which are exactly the terms that
carry no information about which community a document came from. If that happens, $S_{ij}$
collapses toward a constant and stops doing any work.

\paragraph{A measure built on what distinguishes rather than what is shared.}
G.\,P. Morgan's corpora-comparison approach (2017) inverts the question. Rather than asking
how much two bodies of text have in common, it asks, for each token: \emph{if I saw this
word, would I know which of the two corpora it came from?} Tokens are scored by a
transformed normalized odds ratio, so a word that appears constantly in both contributes
nothing, while a word concentrated in one contributes a great deal. The method was
designed precisely to distinguish corpora by differential usage rather than by shared
frequency.

Adapting it here is direct. Treat each community's article text as a corpus. For
communities $i$ and $j$, compute token-level polarization across the two, and define
topical \emph{distance} from the distribution of discriminating terms:
\[
D^{\text{text}}_{ij} \;=\; f\!\left(\{\,\text{polarization}_{w,ij}\,\}_{w \in V}\right),
\qquad
S_{ij} \;=\; 1 - D^{\text{text}}_{ij}
\]
after normalization. The substantive reading is attractive: two communities are topically
proximate when their vocabularies are \emph{hard to tell apart}, not when both happen to
contain the domain's boilerplate.

Morgan's procedure also suggests a preprocessing step we would otherwise have to invent.
Its random-draw stage accumulates evidence about tokens that repeatedly fail to distinguish
randomly constructed subcorpora -- which is a principled way of identifying caffeine-domain
boilerplate before any pairwise distance is computed.

\paragraph{A measurement study, not a commitment.}
We should not adopt this without checking it. Section~\ref{sec:pipeline} provides for a
small measurement study comparing at least three candidates on the same community pairs:

\begin{enumerate}[leftmargin=2em, itemsep=1pt]
  \item corpora-comparison distance in the form above;
  \item a conventional vector similarity over term frequencies, as the baseline the first
        is supposed to improve on;
  \item an embedding-based similarity, if a suitable model is already available.
\end{enumerate}

The validation is human. We select a set of community pairs whose substantive closeness we
can judge from their labels and abstracts -- some obviously related, some obviously not,
some genuinely ambiguous -- and ask which measure orders them the way a reader would. A
measure that cannot separate the obvious cases is not going to be trusted on the ambiguous
ones.

\subsection{Cognitive orientation: two measures, deliberately kept apart}

``Cognitive'' can mean two different things about a piece of research, and in this corpus
they come apart sharply.

\begin{itemize}[leftmargin=2em]
  \item \textbf{Cognitive subject matter}, from abstracts and keywords: does the work study
        attention, vigilance, memory, executive function, reaction time?
  \item \textbf{Cognitive disciplinary provenance}, from journals: is the work published
        where cognitive scientists publish?
\end{itemize}

In era 23, 36.6\% of caffeine articles carry cognitive subject matter \meas{}. Meanwhile
the four core cognitive-architecture venues -- \emph{Cognitive Science}, \emph{Cognitive
Systems Research}, \emph{BICA}, and \emph{Topics in Cognitive Science} -- hold nine nodes
between them in a corpus of 590{,}170 \meas{}.

A vigilance experiment published in \emph{Psychopharmacology} is exactly the case that
makes this worth keeping separate. Its content is cognitive; its provenance is
pharmacological. Collapsing the two into a single ``cognitive contribution'' score would
hide precisely the distinction that might explain why integration succeeds or fails.

We will write $c_i$ for the cognitive orientation of community $i$ -- a proportion between
0 and 1, computed both ways, entering the model as two separate variables rather than an
average.

A caution about the keyword measure: a draft lexicon produced false positives on
\code{inhibit} (enzyme inhibition), \code{memory} (immunological memory), and \code{affect}
used as a verb \meas{}. The lexicon needs pruning and validation against a hand-coded
sample before the measure is trusted, which is one reason the journal-based measure is
worth having alongside it rather than instead of it.

\subsection{Author and venue overlap}

Two communities that share authors, or publish in the same journals, will cite one another
more for reasons that have nothing to do with knowledge integration. Both are measured
directly -- author overlap from the author identifiers in the node list, venue overlap from
the journal profiles -- and both enter the model so that their contribution is not
attributed to the process we care about.

\begin{measure}{How topical proximity is computed}
Three candidates are on the table (\S\ref{sec:proximity}), and the choice matters because
$S_{ij}$ is what distinguishes ``these communities were never going to connect'' from
``these communities should have connected''. The empirical part is a comparison on the same
pairs; the part that needs Sarah is the validation set. Judging which community pairs are
substantively close -- and which are close in a way that \emph{should} produce citation --
is exactly the expertise the citation record does not contain.
\end{measure}

\question{Sarah responded 22 September 2026; her answer is not yet written into the document.}
\begin{decision}{Which cognitive measure carries the primary claim}
Subject matter and disciplinary provenance make different arguments. ``Research on
cognitive questions is not taken up'' is a claim about what is studied; ``work from the
cognitive science community is not taken up'' is a claim about who studies it. The model
can carry both, but the theoretical argument should say which pattern we expect and treat
the other as a competing explanation or a robustness check. Choosing after seeing the
results would not be defensible.
\end{decision}

% =====================================================================
\section{From Networks to Statistical Inference}
\label{sec:inference}
\seccap{Why counting ties is not enough, and the four ideas that make the models
intelligible.}

\subsection{The problem the networks cannot solve}

The networks tell us which communities are connected and how strongly. They do not tell us
whether an observed connection is stronger or weaker than it should have been.

Consider two community pairs, each connected by a single shared work. In one, the citing
community published forty papers and emitted several thousand citations; in the other, it
published two papers and emitted ninety. The same observed tie means something very
different in the two cases. Likewise, some individual works are cited by nearly everyone
regardless of discipline, so a tie running through such a work is weak evidence of anything.

Detecting a \emph{missed} connection therefore requires a model of what citation would
normally look like between communities with comparable opportunities. That is what the rest
of this section builds toward.

\subsection{Four ideas the models rest on}

\paragraph{Exposure.} How many chances did this community have to make this connection? A
community that emits three thousand citations has three thousand opportunities to cite
another literature; one that emits ninety has ninety. Exposure is a known quantity we
compute from the data, and it enters the model as a fixed adjustment rather than something
to be estimated -- we are asserting that twice the opportunity means twice the expected
count, all else equal.

\paragraph{Propensity.} After accounting for opportunity, are some communities simply more
inclined to draw on others, or to be drawn upon? This is not known in advance; it is
estimated. Two communities with identical opportunity can still differ in how outward-facing
they are, and that difference should not be mistaken for the phenomenon of interest.
Operationally, propensity is not measured; it is a pair of varying effects, one attached to
the citing community and one to the cited, that absorb whatever systematic inclination
remains once opportunity and the covariates have been accounted for. A community that cites
widely across the board raises its own effect rather than inflating the estimate for any one
partner.

\paragraph{Work salience.} Some individual works are cited far more than others for reasons
particular to them -- a canonical review, a widely used method. If we ignore this, a single
famous paper shared between two communities inflates their apparent connection. Salience is
likewise estimated rather than measured: it is a varying effect attached to the work itself,
identified because most bridging works are cited by more than one community, so the model
can separate ``this work is cited a great deal everywhere'' from ``this community cites this
work unusually often.'' Section~\ref{sec:modelII} makes both of these concrete -- they are
the terms $\theta_i$, $\psi_j$, and $u_h$ -- and this is the reason the unit of analysis
there is the individual work rather than the community pair.

\paragraph{Partial pooling.} Many communities in this data are small. A community with one
observed tie tells us very little on its own, and treating its observed rate as its true
rate would be overconfident. A multilevel model handles this by estimating each community's
parameter partly from its own data and partly from the field-wide pattern, with the balance
determined by how much data that community actually has. Communities with a lot of data are
barely shrunk toward the average; communities with almost none are shrunk a great deal.

This matters concretely here. Between 38\% and 60\% of participating communities have
exactly one tie \meas{}, which under a simple proportion would give them a connection
strength of 1.0 -- not because their connection is exceptional but because arithmetic
leaves no alternative. Partial pooling is what stops that from happening.

\subsection{Two further ideas used later}

\paragraph{The risk set.} The model has to be told which community pairs \emph{could} have
connected. A pair where the earlier community holds no works eligible to be cited, or the
later community emits no citations at all, is not a case of a connection failing to form;
it is a case where no connection was possible. Those pairs are excluded. The remaining pairs
-- realized and unrealized alike -- are the risk set, and the unrealized ones carry real
information, because a connection that could have formed and did not is exactly what we are
looking for.

\paragraph{Posterior prediction.} These are Bayesian models, so what they produce is not a
single estimate but a distribution over plausible values given the data. That gives us
something a single number cannot: for a particular pair of communities, we can ask what
citation count the model expects and how confident it is, and then ask how far below that
expectation the observed count actually falls. Section~\ref{sec:forgone} builds the
definition of a forgone connection on exactly that.

% =====================================================================
\section{Model I: From Topical Proximity to Common Knowledge}
\label{sec:modelI}
\seccap{Stage 1. Do communities working on related problems end up reading the same
things?}

\subsection{The question and the outcome}

Model~I asks whether topical proximity produces a common knowledge base, and how strongly.

The outcome is the number of works two communities cite in common -- the coupling weight of
\S\ref{sec:networks}. One row of data is one pair of communities in one era pair.

\subsection{Why a negative binomial}

The outcome is a count: it cannot be negative and has no upper bound in principle. The
simplest count model, the Poisson, assumes the variance equals the mean. That assumption
fails here, because community pairs differ enormously in how many chances they had to
overlap, and because some pairs share a burst of works while most share none. The negative
binomial is the same model with one extra parameter allowing the variance to exceed the
mean, which is what this data does.

\subsection{Exposure for coupling}

Two communities can only share works they both had a chance to cite. If community $i$'s
articles cite $R_i$ distinct works and $j$'s cite $R_j$, and works were cited uniformly at
random, the expected overlap would be proportional to $R_iR_j$. But citation is not uniform
-- a small number of works are cited enormously -- so the expected overlap is larger than
the raw product suggests. Writing $d_w$ for the number of times work $w$ is cited,

\begin{equation}
E^{K}_{ijp} \;=\; R_{ip}\,R_{jp}\,\kappa_p,
\qquad
\kappa_p \;=\; \frac{\sum_w d_w^{2}}{\left(\sum_w d_w\right)^{2}},
\end{equation}

where $\kappa_p$ is a single number per era pair carrying the skew of the citation
distribution. If every work were cited exactly once, $\kappa_p$ would be $1/W$ for $W$
works, recovering the uniform case.

\subsection{The model}

Writing $K_{ijp}$ for the coupling count between communities $i$ and $j$ in era pair $p$:

\begin{align}
K_{ijp} &\sim \mathrm{NegBinomial}\!\left(\lambda_{ijp},\, \phi_K\right)\\[2pt]
\log \lambda_{ijp} &= \underbrace{\alpha^{K}_{p}}_{\text{era baseline}}
  + \underbrace{\log E^{K}_{ijp}}_{\text{opportunity}}
  + \underbrace{\alpha_{S,p}\, S_{ijp}}_{\text{topical proximity}}
  + \underbrace{\alpha_{S,\ell}\, S_{ijp}}_{\text{by literature}}
  + \underbrace{\boldsymbol{z}_{ijp}^{\top}\boldsymbol{\gamma}}_{\text{other measures}}
  + \underbrace{\theta^{K}_{i} + \psi^{K}_{j}}_{\text{propensities}}
\end{align}

with the proximity slope allowed to drift across the eleven era pairs,

\begin{equation}
\alpha_{S,p} \sim \mathrm{Normal}\!\left(\bar{\alpha}_S + \gamma_\alpha\,(t_p - \bar{t}),\;
\sigma_\alpha\right),
\qquad
\alpha_{S,\ell} \sim \mathrm{Normal}\!\left(0,\; \sigma_\ell\right).
\end{equation}

\subsection{What each term means, and which question it answers}

\begin{center}
\small
\begin{tabular}{@{}lp{0.72\linewidth}@{}}
\toprule
\textbf{Term} & \textbf{Meaning}\\
\doublerule
$\bar{\alpha}_S$ & How steeply shared reading rises with topical proximity, averaged over
the eleven era pairs. \textbf{This answers RQ1.}\\[3pt]
$\gamma_\alpha$ & Whether that relationship strengthens or weakens over time.
\textbf{Half of RQ4.}\\[3pt]
$\sigma_\ell$ & How much the relationship differs among literatures. If close to zero, the
same rule holds everywhere; if large, some literatures convert proximity into shared
reading far more readily than others. \textbf{Part of RQ7.}\\[3pt]
$\theta^{K}_{i},\psi^{K}_{j}$ & Each community's own tendency to share references, beyond
what opportunity explains.\\[3pt]
$\boldsymbol{z}$ & Author overlap, venue overlap, disciplinary distance, and both cognitive
measures for each community.\\[3pt]
$\phi_K$ & The extra-variability parameter that makes this negative binomial rather than
Poisson.\\
\bottomrule
\end{tabular}
\end{center}

A community pair with high topical proximity and coupling far below what this model expects
is mechanism M2: substantively related, reading different literatures.

% =====================================================================
\section{Model II: From Common Knowledge to Citation Integration}
\label{sec:modelII}
\seccap{Stage 2. Given that two communities read the same things, does citation follow --
and is that weaker for cognitive research?}

\subsection{The question and the outcome}

Model~II asks how much direct citation accompanies a given degree of shared reading, and
whether that conversion is weaker where cognitive research is involved.

The outcome is a citation count. But rather than counting citations between two communities
in aggregate, we count them at the level of the individual work being cited.

\subsection{Why the unit is the individual work}

Some works are cited by almost everyone. If we aggregate to the community-pair level, such a
work inflates every connection it touches and we cannot separate ``these two communities are
genuinely connected'' from ``they both happened to cite a canonical review''. Estimating a
salience term for each work fixes that -- but a salience term needs a work to attach to,
and an aggregated community-pair count has no single work in it. A pair connected by
seventeen shared works would need seventeen salience terms.

So one row of data is one \emph{bridging work} paired with one \emph{citing community}: did
era-$t$ community $i$ cite work $h$, and how often? The cited community $j$ follows from
$h$, since a bridging work belongs to a community in the earlier era by definition. For the
2019--2020 pair the risk set pairs 877 bridging works with every era-23 community that
emits citations -- up to all 394 of them -- giving on the order of 345{,}000 rows, the
overwhelming majority of them zeros \meas{}. Only 151 of those communities appear in a
realized tie; the rest contribute rows that are entirely zero, and those are precisely the
observations RQ10 rests on.

Aggregating the fitted model back over the works belonging to a community recovers the
community-pair quantity, so the question remains one about communities even though the
data are finer.

\subsection{Exposure for citation}

At this resolution the opportunity is simply how many citations the citing community emits:

\begin{equation}
E^{C}_{hip} \;=\; N^{\mathrm{out}}_{ip}.
\end{equation}

A work's own attractiveness does not belong in the exposure term, because that is exactly
what the salience parameter $u_h$ is for.

\subsection{Why the outcome is modeled in two stages}

Citation between two communities involves two substantively distinct things.

First, does any citation relationship form at all? This is a yes-or-no question about
whether the later community is in contact with this piece of the earlier community's output.

Second, given that some citation occurs, how much? This is a question about the strength of
a relationship that already exists.

We model these separately, with one submodel for whether the count clears zero and a second
for how large it is once it has. That separation is a theoretical commitment, not a
convenience. The claim under test is that two communities read the same works and citation
does not follow, which sounds at least as much like a failure of contact as a weakness of an
established connection. Forcing both into one parameter would hide which it is, and the
answer is substantively interesting either way: cognitive orientation might govern whether a
connection forms at all while leaving established connections perfectly normal, or the
reverse. Estimating the two stages separately lets each covariate carry a different
coefficient in each -- which is exactly the comparison RQ9 needs.

\paragraph{What this construction is called, and what it is not.} A two-stage count model of
this kind is a \textbf{hurdle model}. The name is unhelpful but the mechanics are simple: one
process decides whether the outcome clears zero, and if it does, the size is drawn from a
distribution that cannot produce a zero. Every zero in the data therefore comes from the
first stage, and every positive count from the second.

The closely related alternative is a \emph{zero-inflated} model, and the difference matters
because the two are routinely conflated. In a zero-inflated model, zeros arise from two
different sources: some communities are simply incapable of citing this work, and the rest
cite it at some rate that happens to come out as zero. An observed zero is then ambiguous
between the two, and the model has to apportion it. The hurdle form makes no such claim: a
zero means one thing, that contact did not occur.

The hurdle form is the right one here because the risk set has already done the work that
zero inflation would otherwise be doing. Pairs where citation was impossible -- no eligible
works, no emitted citations -- are excluded before estimation (\S\ref{sec:inference}). What
remains are pairs that could have connected, so a zero among them is informative rather than
structural. That is precisely the case the hurdle form describes, and it is the case RQ10
depends on.

\subsection{The model}

For bridging work $h$ belonging to earlier community $j$, cited by later community $i$:

\begin{align}
C_{hip} &\sim \mathrm{Hurdle\text{-}NegBinomial}\!\left(\pi_{hip},\, \mu_{hip},\, \phi_C\right)
\end{align}

where $\pi$ is the probability that any citation occurs and $\mu$ the expected count given
that it does. Both parts take the same form:

\begin{align}
\mathrm{logit}\,\pi_{hip} &= \eta_p + \log E^{C}_{hip}
  + \beta^{(H)}_{K,p}\log(K_{ijp}{+}1)
  + \beta^{(H)}_{Ki}\log(K_{ijp}{+}1)\,c_i
  + \beta^{(H)}_{Kj}\log(K_{ijp}{+}1)\,c_j \nonumber\\
  &\qquad + \boldsymbol{x}_{hijp}^{\top}\boldsymbol{\kappa}
  + u^{(H)}_{h} + \theta^{(H)}_{i} + \psi^{(H)}_{j}\\[4pt]
\log \mu_{hip} &= \alpha^{C}_p + \log E^{C}_{hip}
  + \beta^{(+)}_{K,p}\log(K_{ijp}{+}1)
  + \beta^{(+)}_{Ki}\log(K_{ijp}{+}1)\,c_i
  + \beta^{(+)}_{Kj}\log(K_{ijp}{+}1)\,c_j \nonumber\\
  &\qquad + \boldsymbol{x}_{hijp}^{\top}\boldsymbol{\delta}
  + u^{(+)}_{h} + \theta^{(+)}_{i} + \psi^{(+)}_{j}
\end{align}

\subsection{What each term means, and which question it answers}

\begin{center}
\small
\begin{tabular}{@{}lp{0.70\linewidth}@{}}
\toprule
\textbf{Term} & \textbf{Meaning}\\
\doublerule
$\beta_{K}$ & How much citation accompanies a given degree of shared reading, field-wide.
\textbf{This is RQ2, the baseline everything else is read against.}\\[3pt]
$\beta_{Ki}$ & Whether that conversion is weaker when the \emph{citing} community is
cognitively oriented: cognitive communities not citing work despite shared foundations.
\textbf{Part of RQ9.}\\[3pt]
$\beta_{Kj}$ & Whether it is weaker when the \emph{cited} community is cognitively
oriented: cognitive work not being cited despite common foundations.
\textbf{Part of RQ9.}\\[3pt]
$u_h$ & How citable work $h$ is in its own right, so a canonical paper does not manufacture
apparent connections.\\[3pt]
$\theta_i,\psi_j$ & Each community's tendency to cite and to be cited, beyond opportunity.\\[3pt]
$\boldsymbol{x}$ & Topical proximity, author and venue overlap, disciplinary distance, and
the cited work's vintage -- the last of these being the test of mechanism M1.\\
\bottomrule
\end{tabular}
\end{center}

The two superscripts matter: $\beta^{(H)}$ governs whether a connection forms and
$\beta^{(+)}$ how strong it is once formed, so RQ9 has two answers rather than one.

\subsection{Change over time, and the two community traits}

The baseline conversion and the cognitive terms both drift, so that RQ4 and RQ11 are
estimated rather than inferred by comparing separate fits:

\begin{equation}
\beta^{(\cdot)}_{K,p} \sim \mathrm{Normal}\!\left(\bar{\beta}^{(\cdot)}_{K}
  + \gamma^{(\cdot)}_\beta (t_p - \bar{t}),\; \sigma_\beta\right),
\qquad
\beta^{(\cdot)}_{K\cdot,p} \sim \mathrm{Normal}\!\left(\bar{\beta}^{(\cdot)}_{K\cdot}
  + \gamma^{(\cdot)}_{KC} (t_p - \bar{t}),\; \sigma_{KC}\right).
\end{equation}

RQ3 asks whether being derivative and being generative are related traits. A community in
2019 appears as a citing community in the 2018--2019 pair and as a cited community in the
2019--2020 pair, under the same partition -- so we can ask about both of its tendencies
without ever claiming it persists as the same community. They are given a joint prior:

\begin{equation}
\begin{pmatrix}\theta_c\\ \psi_c\end{pmatrix}
\sim \mathrm{MVNormal}\!\left[
\begin{pmatrix}0\\0\end{pmatrix},\;
\begin{pmatrix}\sigma_\theta^{2} & \rho\,\sigma_\theta\sigma_\psi\\
               \rho\,\sigma_\theta\sigma_\psi & \sigma_\psi^{2}\end{pmatrix}\right],
\end{equation}

where $\rho$ is the correlation RQ3 asks about. Model~I supplies an analogous correlation
for the tendency to share references, which is a different trait and is reported alongside
rather than merged with it.

The model estimates both $\beta_{Ki}$ (cognitive communities citing less than shared foundations predict) and
$\beta_{Kj}$ (cognitive work being cited less than shared foundations predict). These are different claims
about how the field works, and the corpus limit of \S\ref{sec:questions} bears on them
differently. Which pattern we expect, and why, is stated in \S\ref{sec:questions}; it is
registered there rather than here so that it is on record before the model is fit.

\question{Sarah responded 22 September 2026; her answer is not yet written into the document.}
\begin{decision}{Whether Model~I should also move to work resolution}
Model~II sits at work resolution because it needs a salience term. Coupling has the same
dependence -- one work shared by many communities generates related observations -- but its
work-level form needs work--community--community triples rather than work--community pairs,
and we have not counted how many rows that would be. Until we do, Model~I is specified at
the community-pair level with the dependence acknowledged rather than modeled. Counting the
triples is a half-hour job and should be done before the detailed specification.
\end{decision}

% =====================================================================
\section{Identifying Forgone Connections}
\label{sec:forgone}
\seccap{Turning ``this connection should have happened'' into two stated quantities: how far
short, and how sure.}

RQ10 asks which particular community pairs show a connection that should have formed and did
not. Answering it requires a definition that does not depend on how we happen to sort a list,
and that can be argued with -- a reader should be able to disagree with a pair's placement on
the list by disagreeing with a number, not with our taste.

\subsection{What this is not}

Because the language of ``should have happened'' invites the comparison, it is worth saying
plainly what this procedure is not. It is not propensity-score matching, and it is not any
other matched-contrast design.

Matching answers a question of the form: this unit received a treatment; what would have
happened to a comparable unit that did not? It requires a treatment that units either receive
or do not, and it constructs a control group of units with similar propensity to be treated,
so that the difference in outcomes can be read as an effect.

Nothing here has that shape. There is no treatment: cognitive orientation, topical proximity,
and shared reading are traits of communities, not interventions assigned to them. There is no
control group: we are not comparing pair $(i,j)$ to some matched pair $(i',j')$. And the
counterfactual is not another unit's realized outcome; it is a distribution the model
generates for \emph{this} pair, given its own opportunity, its own partners' propensities,
and the salience of the works involved. The comparison is within-pair and model-based,
between what was observed and what the model considered plausible.

The word ``propensity'' does appear in this document, but in an unrelated and more literal
sense: the varying effects $\theta_i$ and $\psi_j$, which capture how outward-facing a
community is. They are parameters of the outcome model, not scores used to build a matched
sample.

\subsection{Why the obvious approach fails}

The natural first move is to rank pairs by how far the observed count falls below the model's
expectation. That does not work, for two reasons.

The gap is measured on the count scale, so it is mechanically larger for pairs with large
expected counts. A pair expected to produce forty citations and producing thirty has a gap of
ten; a pair expected to produce three and producing zero has a gap of three, and is far more
interesting.

The gap also says nothing about confidence. A pair whose expectation is $10 \pm 9$ and a pair
whose expectation is $10 \pm 1$ would rank identically at an observed count of two, though the
second is strong evidence and the first is almost none.

\subsection{Two quantities, doing two different jobs}

The Bayesian machinery gives us a distribution of plausible counts for each pair, and that
distribution supports both questions directly.

\paragraph{How far short.} The expected shortfall is the gap between what the model predicts
and what occurred:

\begin{equation}
D_{ijp} \;=\; \mathbb{E}\!\left[\tilde{C}_{ijp} \;\middle|\; \text{data}\right] \;-\; C_{ijp}.
\end{equation}

This is a count, on the same scale as citations, and it answers ``how much integration is
missing.'' It is the quantity to report when the question is magnitude, and the quantity to
sum over pairs when the question is how much integration the field forgoes in total.

\paragraph{How sure.} The confidence that a shortfall is real, and materially large, is the
posterior probability that the observed count falls below the prediction by more than a
stated margin $\Delta$:

\begin{equation}
F_{ijp} \;=\; \Pr\!\left(\tilde{C}_{ijp} - C_{ijp} > \Delta \;\middle|\; \text{data}\right).
\end{equation}

$F$ is a probability between $0$ and $1$. It is unitless, comparable across pairs of very
different sizes, and it is the quantity to rank on.

The two are not redundant. $D$ can be large while $F$ is low, when the model's prediction is
uncertain -- a pair we cannot say much about. $F$ can be high while $D$ is small, when the
model is confident about a modest shortfall -- a real but minor gap. A pair worth writing
about has both.

\subsection{A worked example}

Consider a citing community $i$ in era~$p$ that publishes 30 articles and emits 2{,}400
citations, and an earlier community $j$ holding 45 distinct works. The two share
$K_{ij} = 22$ commonly cited references, which places them in the upper decile of shared
reading for that era pair, and their topical proximity $S_{ij}$ is high. One is oriented
toward cognitive subject matter and the other is not.

Given that opportunity, that amount of shared reading, and the estimated propensities of both
communities, the model's posterior predictive distribution for citations from $i$ to $j$'s
works centers on 14, with an 80\% interval running from 6 to 24. The observed count is 1.

Then $D_{ij} = 14 - 1 = 13$: thirteen citations' worth of integration did not happen. And
with $\Delta = 3$, $F_{ij} = \Pr(\tilde{C} - 1 > 3) = \Pr(\tilde{C} > 4)$, which given that
predictive distribution is roughly $0.93$. The pair clears both bars and goes on the list.

Now change one thing. Suppose the citing community publishes 2 articles and emits 90
citations. The predictive distribution now centers on $0.5$ and its 80\% interval runs from 0
to 2. An observed count of 0 gives $D_{ij} = 0.5$ and $F_{ij} \approx 0.02$. Nothing is
missing here: the community had almost no opportunity, and silence is what the model expects.
This is the case that a raw proportion or a threshold rule would get wrong, and it is why the
exposure term and the posterior predictive distribution are doing real work rather than
decorating the argument.

\subsection{The reporting rule}

A pair is reported as a forgone connection when four conditions hold together:

\begin{enumerate}[leftmargin=2em, itemsep=1pt]
  \item the two communities share substantial reading, $K_{ijp}$ above a stated quantile of
        the era pair's distribution;
  \item they are substantively proximate, $S_{ij}$ above a stated quantile;
  \item the shortfall is material, $D_{ijp} > \Delta$; and
  \item the model is confident, $F_{ijp}$ above a stated cutoff.
\end{enumerate}

Conditions 1 and 2 are what make the pair a candidate for connection at all; conditions 3 and
4 are what make the absence of connection notable. Reporting $D$ and $F$ together, rather
than a rank, lets a reader see both how much is missing and how sure we are.

\paragraph{From works to pairs.} Model~II estimates at the level of the individual bridging
work, so $\tilde{C}_{ijp}$ is obtained by summing the posterior predictive draws over the
works belonging to $j$ -- draw by draw, so that the resulting distribution carries the
correlation the varying effects induce among works from the same community. Summing point
estimates instead would understate the uncertainty, and $F$ would come out too confident.

\paragraph{Two things the quantity does not claim.} It is not a claim that citation
\emph{would} have occurred under some intervention; the association it departs from is not a
causal effect, for the reasons in \S\ref{sec:causal}. And ``forgone'' means materially less
than expected, not literally zero. Requiring zero would discard the many cases of a
connection that exists but is far thinner than the shared foundations warrant -- which, given
how many ties in this network rest on a single work \meas{}, is likely the more common form
the phenomenon takes.

\question{Sarah responded 22 September 2026; her answer is not yet written into the document.}
\begin{decision}{The materiality margin and the probability threshold}
$\Delta$ and the cutoff on $F$ determine which pairs appear on the list, as do the two
quantile cuts on $K$ and $S$. All four should be set in advance and reported with the
sensitivity analysis described in \S\ref{sec:validation}, showing how the list changes as
they vary. Choosing them after seeing which pairs they produce would make RQ10
unfalsifiable. The substantive question for us is what counts as a material shortfall: is
$\Delta$ a fixed number of citations, or a proportion of the predicted count?
\end{decision}

% =====================================================================
\section{Identifying Major Literatures and Their Lifecycles}
\label{sec:literatures}
\seccap{What a ``literature'' is operationally, how we find them, and how we follow them
through time without claiming communities persist.}

\subsection{What a literature is, operationally}

Everything in this section rests on one definition, so it comes first.

\begin{quote}
\textbf{A literature is a higher-order block: a group of communities that connect to the rest
of the field in the same way.}
\end{quote}

Communities are groups of \emph{articles}, found within a single era. A literature is a group
of \emph{communities}, and it may span eras. It is neither a journal nor a keyword or
label -- those enter later, and for different jobs.

The criterion is structural equivalence in the blockmodel sense, and the phrase ``connect in
the same way'' is doing precise work. Two communities belong to the same block not because
they are densely tied to each other, but because they have the same pattern of ties to
everything else: they draw on the same kinds of earlier work and are drawn on by the same
kinds of later work. Two pharmacokinetic communities in adjacent years may barely cite each
other and still be structurally identical, because the rest of the field treats them
interchangeably. A grouping rule based on density would separate them; a blockmodel does not.
This is why we use one.

\subsection{Discovery, description, and validation are three separate steps}

Conflating these is the standard failure mode of community-labeling work, so we keep them
apart deliberately and give each its own instrument.

\paragraph{Discovery is structural.} The blocks are estimated from the citation network of
communities, and from nothing else. A \textbf{nested degree-corrected stochastic blockmodel}
estimates which groups connect to which, at several levels at once -- communities within
literatures, literatures within broader divisions. \emph{Degree-corrected} means it accounts
for large communities having more connections by virtue of size rather than affinity, which
matters here because community sizes are wildly uneven. The number of groups is chosen by the
data rather than set by us, using a minimum-description-length criterion that balances how
well the grouping fits against how complicated it is: a grouping that needs many blocks to
explain a little more is rejected in favor of a simpler one.

The output is not a single partition but a distribution over them, so each community carries a
probability of belonging to each literature. That uncertainty is worth keeping. A community
sitting ambiguously between two literatures is a substantively interesting object -- it may be
exactly where integration is happening -- and a method that assigned it definitively to one
would conceal that.

\paragraph{Description is textual, and comes second.} Once the blocks exist, we say what each
one is about, from the titles, abstracts, and keywords of its member articles and from the
language-model labels of its member communities. This step is interpretive and it is
downstream: the text never influences which communities are grouped together. Naming a block
``caffeine and cardiovascular risk'' is a claim about what the block contains, and it can be
wrong without the block itself being wrong.

\paragraph{Validation is external.} Journals give an independent partition of the same
communities, derived from publication venue rather than from citation. Comparing the two
answers the second half of RQ7. If the citation-derived literatures line up with journal
boundaries, then what organizes this field is disciplinary affiliation; if they cut across it,
something else is doing the organizing -- which is the more interesting result, and the one
\S\ref{sec:purpose} leads us to expect.

We report agreement with the \textbf{adjusted Rand index}, which measures how much two
partitions of the same objects agree, corrected for the agreement expected by chance alone.
Because a community can publish in several journals and because the blockmodel returns soft
assignments, the headline ARI is computed on hard assignments -- each community to its modal
block and its modal journal -- and accompanied by a soft analogue computed across posterior
draws, so that the reported agreement carries the same uncertainty the partition does. Because
journal identity is not derived from the citation record at any point, the comparison is not
circular. Two caveats belong in the reporting: ARI is sensitive to the number of groups, so
the journal partition should be coarsened to a comparable granularity before comparing; and a
low ARI is a substantive finding here rather than a failure of the method.

\subsection{How this differs from the community detection in \S\ref{sec:paths}}

Two partitioning procedures appear in this document at two different levels, and they should
not be confused.

\begin{center}
\small
\begin{tabular}{@{}lll@{}}
\toprule
 & \textbf{CHAMP (\S\ref{sec:paths})} & \textbf{Nested SBM (here)}\\
\doublerule
Objects partitioned & Articles, within one era & Communities, across eras\\
Produces            & Communities            & Literatures\\
Criterion           & Modularity at a resolution & Description length\\
Chooses its scale by & Sweeping $\gamma$ for a stable range & Model selection, automatic\\
Run                 & Once per era, in Julia & Once over the whole community network\\
\bottomrule
\end{tabular}
\end{center}

CHAMP settles a question modularity maximization cannot settle on its own -- which resolution
parameter to use -- by identifying the range of $\gamma$ over which the partition is stable,
and we apply a pooled $\gamma^{*}$ across eras so that community granularity is comparable
from year to year. It operates on articles. The nested blockmodel then takes those communities
as its nodes. Nothing about CHAMP determines the literatures, and nothing about the
blockmodel revisits the communities.

\subsection{Why not a threshold}

An alternative to a blockmodel is to keep only the strongest ties and see what clusters
remain. We avoid that for a specific and measurable reason: in these networks, between 72\%
and 90\% of ties rest on a single shared work \meas{}. Any threshold either removes almost
the entire network or almost nothing, so the choice of cutoff would determine the answer more
than the data would. This is the general weakness of threshold methods, and it is unusually
severe here. A model that accommodates weak ties with appropriately wide uncertainty is
preferable to one that deletes them and then reports the survivors as though they were the
whole picture.

\subsection{Lifecycles without lineages}

Once communities are assigned to literatures, a literature's trajectory is measured by how
much of the field it occupies in each era: how many articles belong to it, how many
communities, and what share of citation traffic passes through it. Rise, growth, contraction,
and disappearance are movements in those quantities, and because block membership is
probabilistic, each is reported with uncertainty rather than as a count.

This is where the distinction of \S\ref{sec:objects} pays off. None of these measurements
requires knowing that a particular community persisted from one era to the next. A
pharmacokinetic literature can be growing while its constituent communities dissolve and
reform each year, and we can observe the growth without adjudicating the identity claim --
which is the claim Path~A cannot support.

% =====================================================================
\section{Path B: Community Lineages and Persistence}
\label{sec:pathb}
\seccap{The questions that need active-era networks, and the analysis available without
them.}

Section~\ref{sec:paths} set out what active-era networks would allow: following a community
through time, and asking what predicts its survival. The natural model is a
\textbf{discrete-time hazard model} -- for each community in each year, the probability that
it ends that year, given that it has survived so far. Because eras here are single years,
discrete time is the natural form. A \emph{competing risks} version separates two different
endings, dying and being absorbed into another community, which the lifecycle vocabulary
requires and a single-outcome model cannot distinguish.

What that model needs is an identity rule: something that says this year's community is last
year's. Path A cannot supply one. A rule based on citation ties would be built on the same
single-work ties that make thresholding unreliable \meas{}, and an unresolved identity
algorithm should not hold up the rest of the project.

There is, however, a persistence analysis available now, and it needs no identity rule at
all. Individual \emph{works} have observable lifespans: we can see the first year a work is
cited and the last, directly from the edge list. Modeling how long works stay in circulation
-- and whether that depends on the venue, the vintage, or the cognitive orientation of the
communities citing it -- tests something in the review's argument directly. If
pharmacokinetic works stay in circulation longer than cognitive ones, that is evidence about
which knowledge the field retains.

So the division is:

\begin{itemize}[leftmargin=2em]
  \item \textbf{Path A now:} structure, integration, the cognitive case, literature
        lifecycles, and the persistence of cited knowledge.
  \item \textbf{Path B next:} persistence and transformation of the communities themselves.
\end{itemize}

% =====================================================================
\section{Design Matrices}
\label{sec:design}
\seccap{Reference tables: which column implements which quantity.}

The tables below are for looking things up while implementing. Every quantity in them has
been introduced and motivated in the preceding sections; nothing new appears here.

They are included because the gap between a model written in equations and a model a computer
can fit is a table -- one row per observation, one column per quantity -- and most of the
errors that survive a careful specification are errors in building it. Two in particular. The
first is a row that should not exist: a community pair that had no opportunity to connect, or
a bridging work paired with a community that emitted no citations that era. Those rows do not
carry information about a connection failing to form, and including them would inflate the
number of zeros the first stage of Model~II has to explain. The second is a row that should
exist and does not: a pair that could have connected and did not. Those rows are exactly the
evidence RQ10 rests on, and they are the ones a naive construction drops, because they appear
nowhere in the edge list -- they have to be generated from the risk set rather than read off
the data. The pipeline tests in \S\ref{sec:pipeline} check both directions.

The unit of analysis differs between the three tables, and the difference is substantive
rather than clerical. Model~II is at work resolution because that is what identifies the
salience effect $u_h$ (\S\ref{sec:modelII}); Model~I is at pair resolution because coupling is
a property of a pair and has no work-level analogue; the persistence table is at
work--era resolution because it is a survival model, and a survival model needs one row for
each period a unit was at risk. Aggregating the Model~II table over the works belonging to a
community must reproduce the Model~I tie counts exactly, which is one of the consistency
tests.

Each column plays one of five roles. The distinction between \emph{exposure} and
\emph{covariate} is the one to watch: an exposure enters with its coefficient fixed at one on
the log scale, which asserts that twice the opportunity means twice the expected count, while
a covariate's coefficient is estimated from the data. Putting a quantity in the wrong role
changes what the model claims, not just how well it fits.

\begin{center}
\small
\begin{tabular}{@{}lp{0.72\linewidth}@{}}
\toprule
\textbf{Role} & \textbf{Meaning}\\
\doublerule
Index    & Identifies the observational unit -- which work, community, era, or pair this row
           refers to. Not used in estimation.\\
Outcome  & The quantity being explained.\\
Exposure & Known opportunity for the outcome to occur; enters with its coefficient fixed
           rather than estimated.\\
Covariate& A measured characteristic used to explain variation in the outcome.\\
Grouping & An identifier used for partial pooling: rows sharing a value share a varying
           effect.\\
\bottomrule
\end{tabular}
\end{center}

\subsection{Model II: one row per bridging work $\times$ citing community}

A row exists for every combination of a bridging work $h$ and a citing community $i$ in the
risk set -- that is, for every work belonging to an earlier community that a later community
\emph{could} have cited, whether or not it did. The outcome \code{y\_cites} is zero on most
rows, and those zeros are the observations RQ10 is about.

Three grouping columns appear because three varying effects do: \code{work\_id} groups the
salience effect $u_h$, \code{citing\_community} groups the citing propensity $\theta_i$, and
\code{cited\_community} groups the cited propensity $\psi_j$. Note that \code{cited\_community}
is determined by \code{work\_id} rather than being an independent index -- each work belongs to
exactly one earlier community, by the attribution rule the pipeline must fix
(\S\ref{sec:pipeline}). Both cognitive measures appear for both communities, giving four
columns, because \S\ref{sec:networks} keeps subject matter and disciplinary provenance apart
deliberately; the model is fit with one pair and refit with the other.

\begin{center}
\small
\begin{longtable}{@{}>{\raggedright\arraybackslash}p{0.30\linewidth}>{\raggedright\arraybackslash}p{0.13\linewidth}>{\raggedright\arraybackslash}p{0.47\linewidth}@{}}
\toprule
\textbf{Column} & \textbf{Role} & \textbf{Quantity}\\
\doublerule
\endhead
\code{pair\_id}, \code{era\_t}, \code{era\_prior} & index & Era pair\\
\code{work\_id} $h$ & index, grouping & Bridging work; groups the salience effect $u_h$\\
\code{citing\_community} $i$ & index, grouping & Groups $\theta_i$\\
\code{cited\_community} $j$ & index, grouping & Follows from $h$; groups $\psi_j$\\
\midrule
\code{y\_cites} & outcome & Citations from $i$'s articles to $h$; zero where none occurred\\
\midrule
\code{out\_citations\_i} & exposure & $E^{C}$, citations emitted by $i$\\
\midrule
\code{n\_coupling} & covariate & $K_{ij}$, works $i$ and $j$ cite in common\\
\code{topical\_prox} & covariate & $S_{ij}$, from text and venue\\
\code{disc\_distance} & covariate & Disciplinary distance from journal profiles\\
\code{author\_overlap}, \code{venue\_overlap} & covariate & $U_{ij}$, $V_{ij}$\\
\code{cog\_subject\_i}, \code{cog\_subject\_j} & covariate & $c_i$, $c_j$ by subject matter\\
\code{cog\_venue\_i}, \code{cog\_venue\_j} & covariate & $c_i$, $c_j$ by provenance\\
\code{work\_vintage} & covariate & Era of $h$; tests mechanism M1\\
\code{pair\_index} & covariate & 1--11, carrying the time trends\\
\bottomrule
\end{longtable}
\end{center}

\subsection{Model I: one row per community pair}

A row exists for every ordered community pair in the risk set for that era pair. The exposure
is stored as its three components rather than as a single product, so that the degree
correction $\kappa_p$ can be recomputed or varied without rebuilding the frame, and so that
the exposure test in \S\ref{sec:pipeline} has something to check against.

\code{literature\_i} and \code{literature\_j} come from the blockmodel of
\S\ref{sec:literatures}, which means this table cannot be built until that has been fit --
the only ordering dependency between the two models. Because block membership is
probabilistic, the column holds the modal assignment for the primary fit, and the sensitivity
analysis refits across posterior draws of the partition.

\begin{center}
\small
\begin{tabular}{@{}>{\raggedright\arraybackslash}p{0.30\linewidth}>{\raggedright\arraybackslash}p{0.13\linewidth}>{\raggedright\arraybackslash}p{0.47\linewidth}@{}}
\toprule
\textbf{Column} & \textbf{Role} & \textbf{Quantity}\\
\doublerule
\code{pair\_id}, \code{community\_i}, \code{community\_j} & index, grouping & \\
\code{n\_coupling} & outcome & Works cited in common\\
\code{R\_i}, \code{R\_j}, \code{kappa\_p} & exposure & Components of $E^{K}$\\
\code{topical\_prox}, \code{disc\_distance}, \code{author\_overlap}, \code{venue\_overlap},
  \code{cog\_*} & covariate & As above\\
\code{literature\_i}, \code{literature\_j} & grouping & Block membership, for
  $\alpha_{S,\ell}$\\
\bottomrule
\end{tabular}
\end{center}

\subsection{Persistence of cited works}

This is the Path~A persistence analysis of \S\ref{sec:pathb}, and its unit is a
work--era: one row for each era in which a work was still at risk of ceasing to be cited,
running from the era it was first cited until either it stops being cited or the corpus ends.
Works still in circulation at the end are censored rather than treated as having survived
indefinitely.

\code{cog\_share\_citers} is the share of citing communities in that era classified as
cognitively oriented, and it is the column that carries the substantive question: whether the
field retains pharmacokinetic knowledge longer than cognitive knowledge. It is measured per
era rather than once per work, because a work's audience changes.

\begin{center}
\small
\begin{tabular}{@{}>{\raggedright\arraybackslash}p{0.30\linewidth}>{\raggedright\arraybackslash}p{0.13\linewidth}>{\raggedright\arraybackslash}p{0.47\linewidth}@{}}
\toprule
\textbf{Column} & \textbf{Role} & \textbf{Quantity}\\
\doublerule
\code{work\_id}, \code{era} & index & Work and year at risk\\
\code{event}, \code{censored} & outcome & Ceased to be cited; still cited at 2020\\
\code{cog\_share\_citers}, \code{journal}, \code{era\_first\_cited}, \code{n\_communities}
 & covariate & Orientation of citers, venue, vintage, breadth of uptake\\
\bottomrule
\end{tabular}
\end{center}

% =====================================================================
\section{Pipeline Requirements and Tests}
\label{sec:pipeline}
\seccap{What the code must produce, and what must hold if it is producing it correctly.}

\subsection{Order of operations}

The stages run in this order, and the order is not arbitrary -- each depends on the last:

\begin{enumerate}[leftmargin=2em, itemsep=1pt]
  \item Resolve the era files, normalizing three inconsistent directory layouts and two
        partition naming conventions.
  \item Build the era maps, checking that each network's vertices and its partition agree
        in length rather than assuming it.
  \item Build the tie layer: deduplicate the edge list, apply the intersection filter, and
        aggregate once into a table at work resolution. This table is both the source of the
        descriptive networks and the model frame for Model~II, so it is a deliverable rather
        than an intermediate.
  \item Determine which communities participate in any tie. This defines the labeling queue,
        and it is what makes the next stage affordable.
  \item Run Pajek to compute within-community degrees, for participating communities only.
  \item Generate community labels, as a parallel queue against the local language model.
  \item Assemble: join labels, write the networks, partitions, arc tables, and model frames.
\end{enumerate}

Steps 4 and 6 are in that order for a substantive reason. Labels are currently generated
for every community in both eras, but only communities that participate in a tie ever reach
the graph -- which means 67.6\% and 87.9\% of the labeling work is discarded on the two era
pairs examined \meas{}. Computing the ties first removes most of the dominant cost.

\subsection{New quantities the pipeline must produce}

\begin{itemize}[leftmargin=2em, itemsep=1pt]
  \item Exposure: citations emitted per community, distinct references cited per community,
        and the pair-level concentration constant $\kappa_p$.
  \item The risk set, with unrealized pairs represented explicitly rather than omitted.
  \item Topical proximity between communities, from abstracts, keywords, and journals.
  \item \textbf{Journal extraction.} The article-information step currently pulls titles,
        abstracts, and keywords from the Web of Science records but not the source field.
        Journal is needed three times over -- to establish disciplinary breadth, to measure
        provenance separately from subject matter, and to supply the external partition for
        RQ7. Neither route available in the node list is adequate: 83.7\% of node labels
        are bare DOIs, and a DOI prefix names a publisher rather than a journal -- prefix
        \code{10.1016} alone covers 25.7\% of the corpus -- while the remaining 16.3\%
        carry abbreviated journal strings that split one journal across variants, giving
        7{,}099 distinct tokens for 95{,}882 works \meas{}. The Web of Science source
        field is already in the records and is the better instrument.
  \item Author overlap between communities.
  \item Both cognitive measures per community, and vintage per bridging work.
\end{itemize}

\subsection{Defects to repair}

\begin{itemize}[leftmargin=2em, itemsep=1pt]
  \item \textbf{Arbitrary attribution.} When a work is cited by several communities, the
        current code credits it to whichever one happens to sort first in the edge list. This
        deletes 15.9\% of ties on one era pair and 56.2\% on the next \meas{}. Because the
        damage differs by pair, the pairs are not comparable, which is fatal for an analysis
        whose point is comparison across years.
  \item \textbf{Duplicated edges.} Each citation is stored between two and fourteen times,
        non-uniformly \meas{}, so the fix is row deduplication rather than dividing by two.
  \item \textbf{Prompts for small communities.} Communities with fewer than five members
        receive prompts padded with the literal string \code{"NA"}, affecting 11.1\% of one
        era's communities and 6.1\% of the next \meas{}. Since labels are how we recognize
        what a literature is, this sits on the critical path.
\end{itemize}

\subsection{Tests}

The tests below are grouped by what they protect.

\paragraph{Inputs.} Every era resolves to an existing network and partition file; vertex
count equals partition length; each era's vertex set matches the node list exactly (verified
for all twelve eras of the integration window, in both directions \meas{}); after deduplication the edge list has no repeated
sender--target pair and 1{,}238{,}336 rows remain \meas{}.

\paragraph{Determinism.} Shuffling the order of the edge list leaves every output
bit-identical. This is the regression test for the arbitrary-attribution defect and the most
important test in the suite.

\paragraph{Consistency between resolutions.} Aggregating the work-level model frame over the
works belonging to a community reproduces the community-pair tie counts exactly. This is
what guarantees the two resolutions describe the same object.

\paragraph{Conservation.} Fractional weights sum, per era pair, to the number of distinct
bridging works; weights normalized within a community sum to one.

\paragraph{Exposure.} $E^{K}$ equals $R_iR_j\kappa_p$ on a hand-built example with known
reference lists; $\kappa_p$ equals $1/W$ exactly when every work is cited once.

\paragraph{Independence of the proximity measure.} Topical proximity and disciplinary
distance are computed by a routine that cannot see the edge list, enforced by a test rather
than by convention.

\paragraph{Prompts.} No prompt contains the literal string \code{"NA"}; every participating
community receives a label.

\paragraph{Known-answer fixtures.} The 2019--2020 pair reproduces 561 ties over 877 bridging
works, with 151 of 394 citing and 234 of 794 cited communities participating \meas{}; the
2018--2019 pair reproduces 251 ties over 126 bridging works \meas{}. A change to these
numbers is a deliberate specification change and the fixture is updated in the same commit
that causes it.

These fixtures are development fixtures, and they are expected to change once. Eras 1--21
carry partition files named \code{era<N>\_Community.clu} and eras 22--23 carry
\code{era<N>\_testCommunity.clu}. That naming split is not, however, where the granularity
actually breaks. Community counts run 474, 455, then \emph{1{,}126} across eras 17, 18 and
19 \meas{}, and era 19 has fewer vertices than era 18. Size cannot produce that. The cause is
connectivity rather than resolution: era 18's network has a giant component holding 74\% of
its vertices and era 19's largest holds 5.4\% \meas{}, and community counts track component
counts because modularity cannot merge across components. Checked directly, only 5 of era
19's 1{,}112 components are subdivided at all, so 1{,}107 of its 1{,}126 ``communities'' are
connected components \meas{}. That question is not left open: \S\ref{sec:paths} settles
it by regenerating every era's partition with CHAMP at a pooled $\gamma^{*}$, which makes
community granularity comparable across eras by construction. The current partitions are what
the Path~A orchestrator is developed and tested against; the fixture numbers above are
rebuilt from the CHAMP partitions before any substantive result is reported, and no temporal
claim in this document should be made from the current ones.

\subsection{Implementation Status}

Recorded as of 22 September 2026, so that a reader can distinguish code that exists from
intended work. Throughout the document, \meas{} marks a value measured from this corpus
during implementation, rather than one inherited from a source or asserted by argument. The
findings in the following subsection rest almost entirely on such measurements.

\paragraph{Built and passing.} The Path~A foundation is written in R under
\code{scripts/Path\_A/}: Pajek readers that return values rather than writing to the global
environment, era-file resolution across the three directory layouts, era maps, and the tie
layer. Its suite passes 38 checks. The known-answer fixtures reproduce 561 ties over 877
bridging works for 2019--2020 and 251 over 126 for 2018--2019, with 151 of 394 citing and 234
of 794 cited communities participating. These results confirm that the rewritten layer
encodes the same both-periods reading as the original. Shuffling the edge list leaves bridging
works, work-resolution ties, and descent weights bit-identical, so the arbitrary-attribution
defect of \S\ref{sec:pipeline} is repaired and covered by a regression test.

The community-detection stage is written in Julia under
\code{community\_detection\_tools/}: Pajek format I/O, the Leiden stack, and CHAMP, behind
one wrapper module. The Pajek suite passes 182 checks, including every era of the integration
window; the Leiden suite passes 33 and the CHAMP suite 67.

The scale-assessment stage is written and has been run over the whole integration window.
\code{champ\_era\_sweep} sweeps one era and returns its hull coefficients, its selected
resolution, and a row of diagnostics; \code{write\_champ\_sweep} and
\code{write\_champ\_summary} record them; the driver \code{champ\_sweep\_eras.jl} runs eras
12 through 23. Two complete sweeps have been run, one over $\gamma \in [0.1, 1.8]$ and one
over $\gamma \in [0.01, 1.0]$. Their results are the subject of \S\ref{sec:scalefinding}, and
they are the reason the community stage is now blocked on a construction question rather than
on code.

\paragraph{Written, not yet run.} Coupling and the exposure constants
(\S\ref{sec:modelI}), built on sparse incidence matrices so that all three weightings follow
from one multiplication. The hand-built fixtures pass; the corpus implementation remains
untested.

\paragraph{Not started.} Topical proximity, journal extraction, the cognitive measures, the
labeling lane, and the orchestrator itself.

% =====================================================================

\subsection{Open Problem: The Era Networks Carry No Detectable Community Scale}
\label{sec:scalefinding}
\label{sec:starfinding}

\begin{quote}
	\textbf{\itshape The Path~A era networks are almost partitioned before community detection
	runs. They consist largely of disconnected stars, each one article and the references it
	cites, because an arc is kept only when both of its endpoints fall inside the same era ---
	which discards roughly 70\% of the citation record. Two complete CHAMP sweeps show that
	neither of CHAMP's selection criteria finds anything in these graphs: the elbow tracks the
	lower edge of whatever $\gamma$ grid it is given, and no partition holds the modularity
	envelope over a wide interval. The obstacle is the construction, not the algorithm, and it
	must be settled before the community stage can proceed.}
\end{quote}

This subsection records what has been measured, states the construction question the
measurements raise, and lists the diagnostic checks that would answer it. None of it is
resolved. It is written so that the work can be resumed from a cold start.

\blocker{Choose the construction the community stage runs on (decision box at the end of this subsection). Run checks C1--C4 first. Until this is settled, nothing downstream of community detection should be built, because the community is the unit of every model in this document.}

\paragraph{What an era network contains.} Era~12 is the 2009 slice. It holds 18{,}256
vertices: 989 caffeine articles returned by the Web of Science search and 17{,}267 cited
works \meas{}. Of those vertices, 965 emit at least one within-era arc and 17{,}345 receive
one, at a mean of 18.6 arcs per citing article \meas{}. Of the works receiving a citation,
96.8\% are cited exactly once \meas{}. The graph falls into 638 weak components, 494 of which
consist of a single citing article and the works it cites. In every one of those 494
components, each vertex other than the citing article has degree exactly one and emits
nothing \meas{}. Mean degree stays near 2.0 across the integration window, and the same
pattern holds in later eras: era~23 has 1{,}823 emitting articles at a mean of 31.0 arcs each,
with 92.4\% of cited works cited once \meas{}.

\paragraph{Why the networks are this sparse.} The immediate cause is the within-era rule. An
arc survives into an era network only when both endpoints belong to that era, and the twelve
era networks together hold roughly 368{,}000 arcs out of about 1{,}238{,}000 in the
deduplicated edge list \meas{}. Approximately 70\% of the citation record is therefore
invisible to community detection. Whatever else is decided, a construction that restores
structure has to return those arcs to the graph that detection runs on.

The second cause is unresolved, and it matters for how \S\ref{sec:paths} describes Path~A.
The node list does not appear to date cited works by their publication year. The work
labelled \code{critserjk\_1988\_fertilsterilNA} carries year 2009 and era 12; so does the work
with DOI \code{10.1016/S0197-4580(97)00002-X} \meas{}. Every one of the 590{,}170 labels
occurs exactly once, and no cited work appears in more than one era \meas{}. The natural
reading is that a cited work receives the year in which it was \emph{first cited} by an
article in the corpus, not the year it was published. The arithmetic supports that reading:
if arcs were same-year citations in the publication sense, a citing article would have on the
order of one within-era arc, not the 18.6 to 31.0 observed. Under first-citation dating, an
era network is the year's articles together with the references none of them had cited before,
which is a star structure almost by construction.

This has not been confirmed against the edge list, and it should be (check~C1 below). An
earlier version of this subsection attributed the sparsity to the rarity of same-year
citation; that explanation is inconsistent with the observed out-degrees and has been
withdrawn.

\paragraph{What CHAMP does on these networks.} Both sweeps were run directed and unweighted,
with weak components of fewer than four vertices dropped beforehand. That filter removes very
little, because the stars themselves average about twenty vertices: era~12 keeps 18{,}110 of
18{,}256 vertices and 567 of 638 components \meas{}.

\begin{center}
\small
\begin{tabular}{rrrrrr}
\toprule
Era & Giant share & $\gamma^{*}$ & Communities & Components & Widest interval \\
\midrule
12 & 6.1\%  & 0.035 & 568 & 567 & 0.19 \\
13 & 18.4\% & 0.061 & 565 & 560 & 0.14 \\
14 & 11.4\% & 0.061 & 560 & 555 & 0.23 \\
15 & 17.7\% & 0.035 & 534 & 530 & 0.13 \\
16 & 51.8\% & 0.035 & 462 & 449 & 0.22 \\
17 & 72.1\% & 0.035 & 341 & 330 & 0.25 \\
18 & 74.1\% & 0.035 & 310 & 297 & 0.26 \\
19 & 5.4\%  & 0.035 & 948 & 947 & 0.29 \\
20 & 4.8\%  & 0.137 & 908 & 904 & 0.23 \\
21 & 36.2\% & 0.035 & 797 & 786 & 0.32 \\
22 & 58.1\% & 0.035 & 620 & 608 & 0.13 \\
23 & 88.3\% & 0.035 & 229 & 219 & 0.18 \\
\bottomrule
\end{tabular}
\end{center}

\noindent All values are measured \meas{}. Giant share is the largest weak component over all
vertices; $\gamma^{*}$, communities, and components come from the sweep over
$\gamma \in [0.01, 1.0]$ with 40 grid points; the widest interval is the longest stretch of
$\gamma$ over which any single partition holds the upper envelope, computed from the recorded
hull coefficients.

Three readings follow.

\emph{The elbow tracks the grid floor.} In the first sweep, over $[0.1, 1.8]$, seven of twelve
eras selected 0.159, the second point of that grid. In the second sweep, over $[0.01, 1.0]$,
nine of twelve selected 0.035, the second point of \emph{that} grid, and two more selected the
third \meas{}. The selected value fell by roughly the ratio of the two floors. Era~23 moved
from 0.628 to 0.035. Widening the window did not reveal an elbow that the first grid had cut
off; it relocated the answer. The reason is structural: below the first split of the giant
component, modularity prefers one community per component, so the sharpest change of slope on
the envelope is always at the bottom of whatever range is swept.

\emph{No partition is stable across a wide interval.} CHAMP's other standard criterion is the
length of the $\gamma$ interval over which one partition stays on the upper envelope. The
widest such interval in any era is 0.32 out of a swept range of 0.99, and most lie between
0.13 and 0.25 \meas{}. Between 13 and 22 distinct partitions take turns holding the envelope
in each era. In four eras the widest interval runs up against the ceiling at $\gamma = 1.0$.
The elbow criterion is pulled to the floor and the width criterion drifts to the ceiling;
neither is locating a property of the network.

\emph{The selected partition is the component partition.} At the selected resolutions, every
era returns its components plus a handful of splits: era~12 gives 568 communities against 567
kept components, era~23 gives 229 against 219 \meas{}. The gain in modularity over simply
labelling each component ranges from $7.5 \times 10^{-6}$ in era~12 to 0.014 in era~23
\meas{}. Modularity itself exceeds 0.985 in every era, for a reason unrelated to mesoscopic
organization: when communities coincide with components, nearly every arc is internal and the
null-model penalty stays small, so $Q$ approaches one arithmetically. That gain figure is also
resolution-dependent rather than a property of the era. Era~23 showed a gain of 0.47 at the
first sweep's $\gamma^{*} = 0.628$ and 0.014 at the second sweep's $\gamma^{*} = 0.035$
\meas{}.

The one durable result is an ordering. Eras 23, 18, 17, and 22 show the most structure and
eras 19, 20, and 12 the least, in both sweeps, and that ordering follows giant-component
share.

\paragraph{What this changes in the rest of the document.} Three things.

First, the pooled $\gamma^{*}$ described in \S\ref{sec:paths} cannot be estimated from these
networks by either criterion. It is not a matter of choosing between per-era and pooled
estimation; there is no fixed point for either to find.

Second, the size-invariance argument of \S\ref{sec:paths} cannot be tested on them either.
That argument holds only under a condition these networks fail:

\begin{quote}
	Stability of $\gamma^{*}$ across eras is evidence for a common mesoscopic scale \emph{only
	when each era contains enough connected structure for the modularity envelope to
	discriminate among plausible partitions}. Without that condition, invariance in
	$\gamma^{*}$ is not evidence of scale invariance at all; it is evidence of flatness.
\end{quote}

Third, and most consequential, the communities themselves are in question. Every model in
this document takes the community as its unit. If an era's partition is its components, then
most communities are one article and the works it cites, and the objects the models are about
are artifacts of the construction. The tie layer is unaffected --- it reads the edge list and
the era maps, not the era networks, so descent and coupling do not depend on this --- but what
those ties run \emph{between} does.

\paragraph{Diagnostic checks, in the order they should be run.}
\label{sec:checks}
None of these has been done. C1 through C4 are estimable from the edge list and node list in
about an hour and require no community detection.

\begin{enumerate}[leftmargin=2.2em, itemsep=2pt]
	\item \textbf{C1. The dating rule.} \checknote{C1: target year vs.\ earliest citing year, from the edge list.} Join each cited work's recorded year to the earliest
	      year among the articles citing it in the edge list, and count how often the two agree.
	      Near-universal agreement confirms first-citation dating, which would mean
	      \S\ref{sec:paths} describes Path~A incorrectly and that the both-periods reading of
	      descent needs restating.
	\item \textbf{C2. The citation-lag distribution.} \checknote{C2: citation lag; needs publication years parsed from labels.} How many years separate a citing article
	      from the work it cites. This requires true publication years for cited works, which are
	      not in the node list if C1 confirms first-citation dating; they can be parsed from the
	      labels of the non-DOI entries, which carry the form
	      \code{author\_year\_journal}. The lag distribution determines what any window width can
	      recover.
	\item \textbf{C3. The window-width curve.} \checknote{C3: arcs inside a window, giant share and mean degree, for widths 1--12 years.} For widths of one through twelve years, the share
	      of arcs falling inside a window, the giant-component share, mean degree, and the
	      component-size profile. This is the measurement that turns the choice of era width from
	      a guess into a finding. If most of the available connectivity arrives by three years,
	      the cost of a five-year window is not worth paying.
	\item \textbf{C4. The whole-corpus graph.} \checknote{C4: components and degree distribution of the full corpus network.} Component profile and degree distribution of the
	      full 590{,}170-vertex, 1{,}238{,}000-arc network. This is the ceiling any windowing can
	      approach, and it is what a single-partition construction would run on.
	\item \textbf{C5. Directed Leiden at corpus scale.} \checknote{C5: fix the incoming-arc scan, or run undirected, before any corpus-scale detection.} \code{delta\_modularity\_directed\_best!}
	      finds a vertex's incoming arcs by scanning the whole adjacency matrix, which costs
	      $O(nnz)$ per vertex. That is tolerable at 50{,}000 vertices and impossible at 590{,}000.
	      Either pass a precomputed transpose, as the function's own notes anticipate, or run the
	      corpus undirected. This blocks C4 from becoming a detection run.
	\item \textbf{C6. Re-derivation under the repaired tie layer.} \checknote{C6: re-derive descent survival and the single-article community share.} The descent-survival figures of
	      0.54\% and 1.55\%, and the finding that 60.7\% of era-23 communities hold exactly one
	      article, were all measured under the pre-repair tie layer and against the legacy
	      partitions. Both need re-running once a construction is settled.
	\item \textbf{C7. The dominance flag.} \checknote{C7: \code{n\_dominant} over-counts; do not use as a diagnostic until fixed.} \code{champ\_community\_detection} marks more
	      partitions dominant than ever reach the envelope: era~12 flags 34, while only 19 win at
	      any $\gamma$ in the swept range \meas{}. The crossing test and its $10^{-12}$ tolerance
	      admit partitions that never become optimal. No conclusion here depends on this, because
	      the elbow rule is applied to a superset, but \code{n\_dominant} should not be used as a
	      diagnostic until it is fixed.
\end{enumerate}

\question{Open: which construction. Three options, with what each costs in era pairs and answerable questions.}
\begin{decision}{Which construction the community stage uses}
	This is the question that blocks everything downstream, and it should be settled with
	C1--C4 in hand rather than in the abstract. Three options are on the table.

	\emph{(a) Wider eras.} Aggregating to multi-year windows restores every arc whose endpoints
	fall within the window, and structure will appear. The cost is era pairs. The analytic
	window spans twelve years, so non-overlapping five-year windows yield two eras and one
	consecutive pair, where the present design rests on eleven. RQ4 and RQ11 estimate drift
	across pairs, and Model~II's varying effects are indexed by pair; one pair supports none of
	that. Sliding windows would give seven pairs but would make eras overlap, which breaks the
	disjointness the descent relation depends on and moves the construction toward Path~B.

	\emph{(b) One partition over the whole range.} Detect communities once on the full corpus
	network, then treat eras as slices of a fixed community structure. This removes the
	resolution problem entirely: one partition, one $\gamma$, nothing to pool and no scale
	screen. What it gives up is time-varying community structure. Communities become fixed
	objects, which makes the trait questions easier and changes the character of anything about
	communities forming, splitting, or dissolving. Path~B's persistence questions largely
	dissolve with it. It also requires C5 first.

	\emph{(c) One-year eras, with detection on a coupling network.} Keep the present era
	definition, but build the network that detection runs on from shared references rather than
	direct citation: two same-year articles are linked when they cite the same work, with the
	shared works drawn from the full edge list rather than only that era. This keeps all eleven
	pairs and uses every arc. The difficulty is that only articles receive communities this way,
	while descent requires cited works to belong to one; assigning each cited work to the
	community of the article that first cited it would approximate what the current construction
	already does implicitly.

	The choice is not purely technical. Each option changes which of the research questions in
	\S\ref{sec:questions} the design can still answer, and that comparison should be made
	explicitly before any code is written.
\end{decision}

\question{On hold until a construction is chosen.}
\begin{decision}{What the scale screen is, and what failing it means}
	This decision is held in abeyance and becomes live again only if a construction is adopted
	whose networks have structure to screen. Three questions then return, and they constrain one
	another.

	\emph{What is screened on.} Giant-component share is the natural structural proxy, but the
	range of modularity across a sweep measures the quantity of interest directly. A threshold
	specified in advance is preferable to one chosen after observing which eras it excludes. If
	the threshold must instead be empirical, the document should say so explicitly.

	\emph{What the sweep runs on.} The whole disconnected graph, as at present, or the giant
	component alone. These need not agree, because the null model's mass is computed over
	whatever graph is supplied.

	\emph{What happens to an era that fails.} Assign it a $\gamma^{*}$ with a warning, exclude
	it from the size-invariance analysis while still partitioning it at the pooled resolution, or
	fall back to a fixed value. The second appears most defensible, but it should be an explicit
	methodological choice rather than a default.
\end{decision}

\subsection{Model Validation and Predictive Checks}
\label{sec:validation}

The tests above check that the pipeline builds the right table. The checks below check that
the model built on that table is worth believing. Both belong in the pipeline, because both
have to run automatically: a check that is performed once by hand and then forgotten is not a
check.

The organizing idea is that a Bayesian model is a generative story -- it says how the data
could have come about -- and a generative story can be run forward. If we run it forward
and it produces data that look nothing like the caffeine literature, the story is wrong,
whatever its coefficients say. Each of the checks below is a version of running the model
forward and looking at what comes out.

\paragraph{Prior predictive checks, run before the data are seen.} Every parameter in
Models~I and~II carries a prior, and priors on the log or logit scale are notoriously easy to
set badly -- a prior that looks modest as a number can imply absurd counts once exponentiated.
So before fitting, we draw parameters from the priors alone, generate citation counts from
them, and ask whether those counts are of a plausible order. Concretely: do the simulated
counts stay within a few orders of magnitude of what a community emitting two thousand
citations could produce, or does the prior routinely generate a community citing one work ten
million times? Do the simulated hurdle probabilities span the full range rather than piling up
at 0 or 1? A prior that fails this is replaced and the replacement documented. This is a test
of the specification, and it can be run before the data exist.

\paragraph{Posterior predictive checks, run after fitting.} Having fit the model, we draw new
datasets from the posterior and compare them to the observed one. The comparison is on
features we care about rather than on fit indices:

\begin{itemize}[leftmargin=2em, itemsep=1pt]
  \item the proportion of zeros, which the hurdle stage exists to get right;
  \item the full distribution of nonzero counts, particularly the upper tail, since a few
        works are cited very heavily and a model that cannot generate them is misspecified;
  \item the distribution of community out-degree and in-degree in the reconstructed network;
  \item the concentration of ties, the HHI of each community's outgoing weight, which is what
        the tie weightings of \S\ref{sec:weightings} are built to describe.
\end{itemize}

\paragraph{Stratified posterior predictive checks.} An aggregate check can pass while the
model is badly wrong in the region the argument depends on. Since our claims are about
particular kinds of pair, we repeat the checks within strata: by topical proximity (does the
model fit high-proximity pairs as well as low-proximity ones?), by community size, by
cognitive orientation, by literature, and by era. The cognitive stratum matters most. If the
model systematically over-predicts citation for cognitively oriented pairs, then every
cognitively oriented pair will appear to have a shortfall, and RQ10's list will be an artifact
of misfit rather than a finding. This check is the one that stands between the forgone-connection
analysis and a fabricated result, and it should be reported whether or not it is flattering.

\paragraph{Leave-one-out cross-validation.} Several modeling choices in this document are
genuine alternatives rather than settled facts, and LOO -- estimated by Pareto-smoothed
importance sampling, so that it does not require refitting -- gives a common scale on which to
compare them by how well each predicts observations it did not see. The comparisons worth
running are those where a scientific claim rides on the choice:

\begin{itemize}[leftmargin=2em, itemsep=1pt]
  \item hurdle against a plain negative binomial, testing whether the two-stage structure of
        \S\ref{sec:modelII} earns its complexity;
  \item with and without the work-level salience effect $u_h$, testing whether individual
        works really do differ in the way \S\ref{sec:inference} assumes;
  \item cognitive orientation by subject matter against orientation by disciplinary
        provenance, which is a substantive question about which measure captures the
        phenomenon and not merely a robustness check;
  \item with and without the coupling term, and with coupling entered linearly against
        entered as a spline, testing whether the relationship in RQ9 is monotone;
  \item literature-level varying intercepts against none, testing whether the blocks of
        \S\ref{sec:literatures} explain anything the covariates do not.
\end{itemize}

Two cautions. Any observation with a Pareto $\hat{k}$ above $0.7$ is refit exactly rather than
approximated, and the count of such observations is reported -- with heavy-tailed citation
counts there will be some. And LOO here is ordinary leave-one-out, which asks how well the
model predicts a held-out \emph{row}, not a held-out era; it is not a test of forecasting
into a future era, and we do not report it as one.

\paragraph{Sensitivity analysis.} The forgone-connection list depends on four stated
constants: $\Delta$, the cutoff on $F$, and the two quantile cuts on $K$ and $S$
(\S\ref{sec:forgone}). We report how the list changes as each varies over a stated range,
including which pairs enter and leave, and we report the same for the choice between the two
cognitive measures. A result that survives only one setting of these constants is reported as
such.

\paragraph{What the pipeline must do with all of this.} Each check above is a function that
returns a pass, a fail, or a number, and the fitting stage writes them to a report alongside
the posterior. The prior predictive check runs before fitting and blocks it on failure. The
posterior and stratified checks and the LOO comparisons run after fitting and are recorded
whether or not they pass, so that a degradation between one fit and the next is visible in the
diff rather than discovered later.

\end{document}
